
\Data\ID{1003019001}
\Updated{2020-07-15 03:07:12}
\Level{A}
\SubArea{Exponential functions}
\LANGen{
\Question{
Which of the following functions is not an exponential function?}
{\( f(x)=x^3 \)}{\( g(x)= \mathrm{e}^{-3x} \)}{\( h(x)= 5^{\frac x3} \)}{\( i(x)= \left(\frac53\right)^x \)}{}{}}
\LANGpl{
\Question{
Która z poniższych funkcji  nie jest funkcją wykładniczą?}
{\( f(x)=x^3 \)}{\( g(x)= \mathrm{e}^{-3x} \)}{\( h(x)= 5^{\frac x3} \)}{\( i(x)= \left(\frac53\right)^x \)}{}{}}
\LANGcs{
\Question{
Která z následujících funkcí není exponenciální?}
{\( f(x)=x^3 \)}{\( g(x)= \mathrm{e}^{-3x} \)}{\( h(x)= 5^{\frac x3} \)}{\( i(x)= \left(\frac53\right)^x \)}{}{}}
\LANGsk{
\Question{
Ktorá z nasledujúcich funkcií nie je exponenciálna?}
{\( f(x)=x^3 \)}{\( g(x)= \mathrm{e}^{-3x} \)}{\( h(x)= 5^{\frac x3} \)}{\( i(x)= \left(\frac53\right)^x \)}{}{}}
\LANGes{
\Question{
¿Cuál de las siguientes funciones no es una función exponencial?}
{\( f(x)=x^3 \)}{\( g(x)= \mathrm{e}^{-3x} \)}{\( h(x)= 5^{\frac x3} \)}{\( i(x)= \left(\frac53\right)^x \)}{}{}}
\End

\Data\ID{1003022401}
\Updated{2021-04-11 09:04:34}
\Level{A}
\SubArea{Exponential functions}
\LANGen{
\Question{
Which of the equations is the equation of the horizontal asymptote of \( f(x)=2^x+3 \)?}
{\( y=3 \)}{\( y=-3 \)}{\( x=3 \)}{\( x=-3 \)}{}{}}
\LANGpl{
\Question{
Które  z podanych równań jest równaniem asymptoty poziomej \( f(x)=2^x+3 \)?}
{\( y=3 \)}{\( y=-3 \)}{\( x=3 \)}{\( x=-3 \)}{}{}}
\LANGcs{
\Question{
Která rovnice je rovnicí vodorovné asymptoty funkce \( f(x)=2^x+3 \)?}
{\( y=3 \)}{\( y=-3 \)}{\( x=3 \)}{\( x=-3 \)}{}{}}
\LANGsk{
\Question{
Ktorá rovnica je rovnicou vodorovnej asymptoty funkcie \( f(x)=2^x+3 \)?}
{\( y=3 \)}{\( y=-3 \)}{\( x=3 \)}{\( x=-3 \)}{}{}}
\LANGes{
\Question{
¿Cuál de las ecuaciones siguientes es la asíntota horizontal de \( f(x)=2^x+3 \)?}
{\( y=3 \)}{\( y=-3 \)}{\( x=3 \)}{\( x=-3 \)}{}{}}
\End

\Data\ID{1103019002}
\Updated{2021-04-11 09:04:49}
\Level{A}
\SubArea{Exponential functions}
\LANGen{
\Question{
One of the given graphs is a graph of an exponential function. Choose this graph.}
{}{}{}{}{}{}}
\LANGpl{
\Question{
Jeden z podanych wykresów jest wykresem funkcji wykładniczej. Który?}
{}{}{}{}{}{}}
\LANGcs{
\Question{
Který z následujících grafů je grafem exponenciální funkce?}
{}{}{}{}{}{}}
\LANGsk{
\Question{
Ktorý z nasledujúcich grafov je graf exponenciálnej funkcie?}
{}{}{}{}{}{}}
\LANGes{
\Question{
¿Cuál de las gráficas es una gráfica de la función exponencial?}
{}{}{}{}{}{}}

\IMGanswer{1103019002a_0.pdf}
\IMGanswer{1103019002b_0.pdf}
\IMGanswer{1103019002c_0.pdf}
\IMGanswer{1103019002d_0.pdf}\End

\Data\ID{1103022402}
\Updated{2020-07-15 03:07:12}
\Level{A}
\SubArea{Exponential functions}
\IMG{1103022402_0.pdf}
\LANGen{
\Question{
Which of the equations is the equation of the horizontal asymptote of \( f(x) = \left(\frac12\right)^x-2 \)? Function \( f \) is graphed in the figure.}
{\( y=-2 \)}{\( y=-3 \)}{\( x=-2 \)}{\( x=-3 \)}{}{}}
\LANGpl{
\Question{
Które z podanych równań jest równaniem asymptoty poziomej funkcji  \( f(x) = \left(\frac12\right)^x-2 \)? Wykres funkcji  \( f \) przedstawiono na rysunku.}
{\( y=-2 \)}{\( y=-3 \)}{\( x=-2 \)}{\( x=-3 \)}{}{}}
\LANGcs{
\Question{
Na obrázku je graf funkce \( f(x) = \left(\frac12\right)^x-2 \). Která rovnice je rovnicí vodorovné asymptoty dané funkce \( f \)?}
{\( y=-2 \)}{\( y=-3 \)}{\( x=-2 \)}{\( x=-3 \)}{}{}}
\LANGsk{
\Question{
Na obrázku je graf funkcie \( f(x) = \left(\frac12\right)^x-2 \). Ktorá rovnica je rovnicou vodorovnej asymptoty danej funkcie \( f \)?}
{\( y=-2 \)}{\( y=-3 \)}{\( x=-2 \)}{\( x=-3 \)}{}{}}
\LANGes{
\Question{
En la imagen está la función \( f(x) = \left(\frac12\right)^x-2 \). ¿Cuál de las ecuaciones es la ecuación de la asíntota horizontal de la función \( f \)?}
{\( y=-2 \)}{\( y=-3 \)}{\( x=-2 \)}{\( x=-3 \)}{}{}}
\End

\Data\ID{1103022403}
\Updated{2020-07-15 03:07:10}
\Level{A}
\SubArea{Exponential functions}
\IMG{1103022403_0.pdf}
\LANGen{
\Question{
The graph of exponential function \( f(x) = \left(\frac13\right)^x \) is below. Use the graph of \( f \) to identify which of the following graphs is the graph of \( g(x) = \left(\frac13\right)^x + 2\).}
{}{}{}{}{}{}}
\LANGpl{
\Question{
Wykres funkcji wykładniczej \( f(x) = \left(\frac13\right)^x \) przedstawiono poniżej.  Używając wykresu funkcji  \( f \) określ, który z podanych wykresów jest wykresem funkcji \( g(x) = \left(\frac13\right)^x + 2\).}
{}{}{}{}{}{}}
\LANGcs{
\Question{
Na obrázku je graf funkce \( f(x) = \left(\frac13\right)^x \). Pomocí grafu funkce \( f \) identifikujte, který z následujících grafů je grafem funkce \( g(x) = \left(\frac13\right)^x + 2\).}
{}{}{}{}{}{}}
\LANGsk{
\Question{
Na obrázku je graf funkcie \( f(x) = \left(\frac13\right)^x \). Pomocou grafu funkcie \( f \) identifikujte, ktorý z nasledujúcich grafov je graf funkcie \( g(x) = \left(\frac13\right)^x + 2\).}
{}{}{}{}{}{}}
\LANGes{
\Question{
En la imagen está la gráfica de la función \( f(x) = \left(\frac13\right)^x \). Usando la gráfica de la función \( f \) identifica cuál de las siguientes gráficas es la gráfica de \( g(x) = \left(\frac13\right)^x + 2\).}
{}{}{}{}{}{}}

\IMGanswer{1103022403a_0.pdf}
\IMGanswer{1103022403b_0.pdf}
\IMGanswer{1103022403c_0.pdf}
\IMGanswer{1103022403d_0.pdf}\End

\Data\ID{1103022404}
\Updated{2020-07-15 03:07:10}
\Level{A}
\SubArea{Exponential functions}
\IMG{1103022404.pdf}
\LANGen{
\Question{
The graph of exponential function \( f(x) = 6^x \) is below. Use the graph of \( f \) to identify which of the following graphs is the graph of \( g(x) = 6^{x-3} \).}
{}{}{}{}{}{}}
\LANGpl{
\Question{
Wykres funkcji wykładniczej \( f(x) = 6^x \) przedstawiono poniżej. Używając wykresu funkcji \( f \) określ, który z poniższych wykresów jest wykresem funkcji \( g(x) = 6^{x-3} \).}
{}{}{}{}{}{}}
\LANGcs{
\Question{
Na obrázku je graf exponenciální funkce \( f(x) = 6^x \). Pomocí grafu funkce \( f \) identifikujte, který z následujících grafů je graf funkce \( g(x) = 6^{x-3} \).}
{}{}{}{}{}{}}
\LANGsk{
\Question{
Na obrázku je graf exponenciálnej funkcie \( f(x) = 6^x \). Pomocou grafu funkcie \( f \) identifikujte, ktorý z nasledujúcich grafov je graf funkcie \( g(x) = 6^{x-3} \).}
{}{}{}{}{}{}}
\LANGes{
\Question{
En la imagen está la gráfica de la función \( f(x) = 6^x \). Usa la gráfica de la función \( f \) para identificar cuál de las siguientes gráficas es la representación de \( g(x) = 6^{x-3} \).}
{}{}{}{}{}{}}

\IMGanswer{1103022404a_0.pdf}
\IMGanswer{1103022404b_0.pdf}
\IMGanswer{1103022404c.pdf}
\IMGanswer{1103022404d_0.pdf}\End

\Data\ID{1103022405}
\Updated{2020-07-15 03:07:10}
\Level{A}
\SubArea{Exponential functions}
\IMG{1103022405_0.pdf}
\LANGen{
\Question{
The graph of exponential function \( f(x) = 3^{x-2} \) is below. Use the graph of \( f \) to identify which of the following graphs is the graph of \( g(x) = -3^{x-2} \).}
{}{}{}{}{}{}}
\LANGpl{
\Question{
Wykres funkcji wykładniczej \( f(x) = 3^{x-2} \)  przedstawiono poniżej. Korzystając  z wykresu funkcji  \( f \) określ, który z poniższych wykresów jest wykresem funkcji \( g(x) = -3^{x-2} \).}
{}{}{}{}{}{}}
\LANGcs{
\Question{
Na obrázku je graf funkce \( f(x) = 3^{x-2} \). Pomocí grafu funkce \( f \) identifikujte, který z následujících grafů je graf funkce \( g(x) = -3^{x-2} \).}
{}{}{}{}{}{}}
\LANGsk{
\Question{
Na obrázku je graf funkcie \( f(x) = 3^{x-2} \). Pomocou grafu funkcie \( f \) identifikujte, ktorý z nasledujúcich grafov je graf funkcie \( g(x) = -3^{x-2} \).}
{}{}{}{}{}{}}
\LANGes{
\Question{
En la imagen está la representación gráfica de la función \( f(x) = 3^{x-2} \). Usa la gráfica de la \( f \) para identificar cuál de las gráficas es la representación de \( g(x) = -3^{x-2} \).}
{}{}{}{}{}{}}

\IMGanswer{1103022405a_0.pdf}
\IMGanswer{1103022405b_0.pdf}
\IMGanswer{1103022405c_0.pdf}
\IMGanswer{1103022405d_0.pdf}\End

\Data\ID{1103022406}
\Updated{2020-07-15 03:07:10}
\Level{A}
\SubArea{Exponential functions}
\IMG{1103022406.pdf}
\LANGen{
\Question{
The graph of exponential function \( f(x) = 3^{x-2} \)  is below.  Use the graph of \( f \) to identify which of the following graphs is the graph of \( g(x) = 3^{-(x-2)} \).}
{}{}{}{}{}{}}
\LANGpl{
\Question{
Wykres funkcji wykładniczej \( f(x) = 3^{x-2} \)  przedstawiono poniżej.  Korzystając z wykresu funkcji \( f \) określ, który z podanych wykresów jest wykresem funkcji  \( g(x) = 3^{-(x-2)} \).}
{}{}{}{}{}{}}
\LANGcs{
\Question{
Na obrázku je graf exponenciální funkce \( f(x) = 3^{x-2} \). Pomocí grafu funkce \( f \) identifikujte, který z následujících grafů je graf funkce \( g(x) = 3^{-(x-2)} \).}
{}{}{}{}{}{}}
\LANGsk{
\Question{
Na obrázku je graf exponenciálnej funkcie \( f(x) = 3^{x-2} \). Pomocou grafu funkcie \( f \) identifikujte, ktorý z nasledujúcich grafov je graf funkcie \( g(x) = 3^{-(x-2)} \).}
{}{}{}{}{}{}}
\LANGes{
\Question{
En la imagen está la gráfica de la función \( f(x) = 3^{x-2} \).  Usa la gráfica de la función \( f \) para identificar cuál de las siguientes gráficas es la representación de la función \( g(x) = 3^{-(x-2)} \).}
{}{}{}{}{}{}}

\IMGanswer{1103022406a.pdf}
\IMGanswer{1103022406b.pdf}
\IMGanswer{1103022406c.pdf}
\IMGanswer{1103022406d.pdf}\End

\Data\ID{1103022407}
\Updated{2020-11-12 01:11:24}
\Level{A}
\SubArea{Exponential functions}
\LANGen{
\Question{
Which of the following graphs is the graph of function \( f(x) = 4-\left(\frac13\right)^{(x-2)}\)?}
{}{}{}{}{}{}}
\LANGpl{
\Question{
Który z podanych wykresów jest wykresem funkcji \( f(x) = 4-\left(\frac13\right)^{(x-2)}\)?}
{}{}{}{}{}{}}
\LANGcs{
\Question{
Který z následujících grafů je graf funkce \( f(x) = 4-\left(\frac13\right)^{(x-2)}\)?}
{}{}{}{}{}{}}
\LANGsk{
\Question{
Ktorý z nasledujúcich grafov je graf funkcie \( f(x) = 4-\left(\frac13\right)^{(x-2)}\)?}
{}{}{}{}{}{}}
\LANGes{
\Question{
¿Cuál de las siguientes gráficas es la representación de la función \( f(x) = 4-\left(\frac13\right)^{(x-2)}\)?}
{}{}{}{}{}{}}

\IMGanswer{1103022407a.pdf}
\IMGanswer{1103022407b_0.pdf}
\IMGanswer{1103022407c.pdf}
\IMGanswer{1103022407d.pdf}\End

\Data\ID{1103022408}
\Updated{2020-11-12 01:11:55}
\Level{A}
\SubArea{Exponential functions}
\LANGen{
\Question{
Which of the given graphs is the graph of function \( f(x)=5\cdot 0.2^x\)?}
{}{}{}{}{}{}}
\LANGpl{
\Question{
Który z poniższych wykresów jest  wykresem funkcji \( f(x)=5\cdot 0{,}2^x\)?}
{}{}{}{}{}{}}
\LANGcs{
\Question{
Který z následujících grafů je graf funkce \( f(x)=5\cdot 0{,}2^x\)?}
{}{}{}{}{}{}}
\LANGsk{
\Question{
Ktorý z nasledujúcich grafov je graf funkcie \( f(x)=5\cdot 0{,}2^x\)?}
{}{}{}{}{}{}}
\LANGes{
\Question{
¿Cuál de las gráficas dadas es la representación de la función \( f(x)=5\cdot 0.2^x\)?}
{}{}{}{}{}{}}

\IMGanswer{1103022408a_0.pdf}
\IMGanswer{1103022408b_0.pdf}
\IMGanswer{1103022408c_0.pdf}
\IMGanswer{1103022408d_0.pdf}\End

\Data\ID{1103037401}
\Updated{2020-11-16 11:11:48}
\Level{A}
\SubArea{Exponential functions}
\IMG{1103037401.pdf}
\LANGen{
\Question{
The graph of the function \(3^x-9\) is below. What is the domain of \( f \)?}
{\( (-\infty;\infty ) \)}{\( (-9;\infty ) \)}{\( (-\infty;9) \)}{\( [-9;\infty ) \)}{}{}}
\LANGpl{
\Question{
Wykres funkcji \(3^x-9\) przedstawiono poniżej. Jaka jest dziedzina funkcji \( f \)?}
{\( (-\infty;\infty ) \)}{\( (-9;\infty ) \)}{\( (-\infty;9) \)}{\( [-9;\infty ) \)}{}{}}
\LANGcs{
\Question{
Na obrázku je graf funkce \(3^x-9\). Určete definiční obor funkce \( f \).}
{\( (-\infty;\infty ) \)}{\( (-9;\infty ) \)}{\( (-\infty;9) \)}{\( \langle -9;\infty ) \)}{}{}}
\LANGsk{
\Question{
Na obrázku je graf funkcie \(3^x-9\). Určte definičný obor funkcie \( f \).}
{\( (-\infty;\infty ) \)}{\( (-9;\infty ) \)}{\( (-\infty;9) \)}{\( \langle -9;\infty ) \)}{}{}}
\LANGes{
\Question{
En la imagen está la función \(3^x-9\). Determina el dominio de esta función \(f\).}
{\( (-\infty;\infty ) \)}{\( (-9;\infty ) \)}{\( (-\infty;9) \)}{\( [-9;\infty ) \)}{}{}}
\End

\Data\ID{1103037402}
\Updated{2020-11-16 11:11:12}
\Level{A}
\SubArea{Exponential functions}
\IMG{1103037402.pdf}
\LANGen{
\Question{
The graph of the function \(-0.3^x+4\) is below. What is the range of  \( f \)?}
{\( (-\infty;4) \)}{\( (4;\infty) \)}{\( (-\infty;4] \)}{\( (-\infty;\infty) \)}{}{}}
\LANGpl{
\Question{
Wykres funkcji \(-0{,}3^x+4\) przedstawiono poniżej. Jaki jest zakres funkcji \( f \)?}
{\( (-\infty;4) \)}{\( (4;\infty) \)}{\( (-\infty;4] \)}{\( (-\infty;\infty) \)}{}{}}
\LANGcs{
\Question{
Na obrázku je graf funkce \(-0{,}3^x+4\). Určete obor hodnot funkce \( f \).}
{\( (-\infty;4) \)}{\( (4;\infty) \)}{\( (-\infty;4 \rangle\)}{\( (-\infty;\infty) \)}{}{}}
\LANGsk{
\Question{
Na obrázku je graf funkcie \(-0{,}3^x+4\). Určte obor hodnôt funkcie \( f \).}
{\( (-\infty;4) \)}{\( (4;\infty) \)}{\( (-\infty;4 \rangle\)}{\( (-\infty;\infty) \)}{}{}}
\LANGes{
\Question{
En la imagen está la función \(-0.3^x+4\). Determina el rango de esta función \(f\).}
{\( (-\infty;4) \)}{\( (4;\infty) \)}{\( (-\infty;4] \)}{\( (-\infty;\infty) \)}{}{}}
\End

\Data\ID{9000003701}
\Updated{2021-04-11 09:04:14}
\Level{A}
\SubArea{Exponential functions}
\IMG{000037_01-figure0.pdf}
\LANGen{
\Question{
Identify a possible analytic expression for the exponential function graphed in the
picture.}
{\(y = 2^{x-1} - 2\)}{\(y = 2^{x+1} - 2\)}{\(y = 2^{x+1} + 2\)}{\(y = 2^{x-1} + 2\)}{}{}}
\LANGpl{
\Question{
Określ możliwy wzór funkcji wykładniczej, której wykres przedstawiono poniżej.}
{\(y = 2^{x-1} - 2\)}{\(y = 2^{x+1} - 2\)}{\(y = 2^{x+1} + 2\)}{\(y = 2^{x-1} + 2\)}{}{}}
\LANGcs{
\Question{
Určete předpis funkce, jejíž graf je znázorněn na obrázku.}
{\(y = 2^{x-1} - 2\)}{\(y = 2^{x+1} - 2\)}{\(y = 2^{x+1} + 2\)}{\(y = 2^{x-1} + 2\)}{}{}}
\LANGsk{
\Question{
Určte predpis funkcie, ktorej graf je znázornený na obrázku.}
{\(y = 2^{x-1} - 2\)}{\(y = 2^{x+1} - 2\)}{\(y = 2^{x+1} + 2\)}{\(y = 2^{x-1} + 2\)}{}{}}
\LANGes{
\Question{
Identifica la expresión analítica de la función exponencial representada por la gráfica.}
{\(y = 2^{x-1} - 2\)}{\(y = 2^{x+1} - 2\)}{\(y = 2^{x+1} + 2\)}{\(y = 2^{x-1} + 2\)}{}{}}
\End

\Data\ID{9000003702}
\Updated{2021-04-11 09:04:38}
\Level{A}
\SubArea{Exponential functions}
\LANGen{
\Question{
In the following list identify a function with graph through the points
\([3;0]\) and
\([5;3]\).}
{\(y = \left (\frac{1}
{2}\right )^{3-x} - 1\)}{\(y = \left (\frac{1}
{2}\right )^{3-x} + 1\)}{\(y = 1 -\left (\frac{1}
{2}\right )^{x-3}\)}{\(y = \left (\frac{1}
{2}\right )^{x-3} + 1\)}{\(y = 1 -\left (\frac{1}
{2}\right )^{x+3}\)}{\(y = \left (\frac{1}
{2}\right )^{x-3} - 1\)}}
\LANGpl{
\Question{
Z poniższych funkcji wybierz tę, której wykres przechodzi przez punkty
\([3;0]\) i
\([5;3]\).}
{\(y = \left (\frac{1}
{2}\right )^{3-x} - 1\)}{\(y = \left (\frac{1}
{2}\right )^{3-x} + 1\)}{\(y = 1 -\left (\frac{1}
{2}\right )^{x-3}\)}{\(y = \left (\frac{1}
{2}\right )^{x-3} + 1\)}{\(y = 1 -\left (\frac{1}
{2}\right )^{x+3}\)}{\(y = \left (\frac{1}
{2}\right )^{x-3} - 1\)}}
\LANGcs{
\Question{
Funkce, jejíž graf prochází body
\([3;0]\),
\([5;3]\), má
předpis:}
{\(y = \left (\frac{1}
{2}\right )^{3-x} - 1\)}{\(y = \left (\frac{1}
{2}\right )^{3-x} + 1\)}{\(y = 1 -\left (\frac{1}
{2}\right )^{x-3}\)}{\(y = \left (\frac{1}
{2}\right )^{x-3} + 1\)}{\(y = 1 -\left (\frac{1}
{2}\right )^{x+3}\)}{\(y = \left (\frac{1}
{2}\right )^{x-3} - 1\)}}
\LANGsk{
\Question{
Funkcia, ktorej graf prechádza bodmi
\([3;0]\),
\([5;3]\), má
predpis:}
{\(y = \left (\frac{1}
{2}\right )^{3-x} - 1\)}{\(y = \left (\frac{1}
{2}\right )^{3-x} + 1\)}{\(y = 1 -\left (\frac{1}
{2}\right )^{x-3}\)}{\(y = \left (\frac{1}
{2}\right )^{x-3} + 1\)}{\(y = 1 -\left (\frac{1}
{2}\right )^{x+3}\)}{\(y = \left (\frac{1}
{2}\right )^{x-3} - 1\)}}
\LANGes{
\Question{
Identifica la función cuya gráfica que pasa por los puntos 
\([3;0]\) y
\([5;3]\).}
{\(y = \left (\frac{1}
{2}\right )^{3-x} - 1\)}{\(y = \left (\frac{1}
{2}\right )^{3-x} + 1\)}{\(y = 1 -\left (\frac{1}
{2}\right )^{x-3}\)}{\(y = \left (\frac{1}
{2}\right )^{x-3} + 1\)}{\(y = 1 -\left (\frac{1}
{2}\right )^{x+3}\)}{\(y = \left (\frac{1}
{2}\right )^{x-3} - 1\)}}
\End

\Data\ID{9000003703}
\Updated{2020-07-15 03:07:34}
\Level{A}
\SubArea{Exponential functions}
\LANGen{
\Question{
In the following list identify a point which is not on the graph of the function
\(f(x) = 3 -\left (\frac{1}
{3}\right )^{x}\).}
{\(C = [-2;6]\)}{\(A = [-1;0]\)}{\(B = \left [1; \frac{8} 
{3}\right ]\)}{\(D = [0;2]\)}{\(E = [-3;-24]\)}{\(F = \left [2; \frac{26} 
 {9} \right ]\)}}
\LANGpl{
\Question{
Z poniższych punktów wybierz ten, który nie należy do wykresu funkcji
\(f(x)= 3 -\left (\frac{1}
{3}\right )^{x}\).}
{\(C = [-2;6]\)}{\(A = [-1;0]\)}{\(B = \left [1; \frac{8} 
{3}\right ]\)}{\(D = [0;2]\)}{\(E = [-3;-24]\)}{\(F = \left [2; \frac{26} 
 {9} \right ]\)}}
\LANGcs{
\Question{
Kterým bodem neprochází graf funkce
\(f(x) = 3 -\left (\frac{1}
{3}\right )^{x}\)?}
{\(C = [-2;6]\)}{\(A = [-1;0]\)}{\(B = \left [1; \frac{8} 
{3}\right ]\)}{\(D = [0;2]\)}{\(E = [-3;-24]\)}{\(F = \left [2; \frac{26} 
 {9} \right ]\)}}
\LANGsk{
\Question{
Ktorým bodom neprechádza graf funkcie
\(f(x) = 3 -\left (\frac{1}
{3}\right )^{x}\)?}
{\(C = [-2;6]\)}{\(A = [-1;0]\)}{\(B = \left [1; \frac{8} 
{3}\right ]\)}{\(D = [0;2]\)}{\(E = [-3;-24]\)}{\(F = \left [2; \frac{26} 
 {9} \right ]\)}}
\LANGes{
\Question{
Entre los siguientes puntos elige el cual no pasa por la gráfica de la función. 
\(f(x) = 3 -\left (\frac{1}
{3}\right )^{x}\).}
{\(C = [-2;6]\)}{\(A = [-1;0]\)}{\(B = \left [1; \frac{8} 
{3}\right ]\)}{\(D = [0;2]\)}{\(E = [-3;-24]\)}{\(F = \left [2; \frac{26} 
 {9} \right ]\)}}
\End

\Data\ID{9100003606}
\Updated{2020-07-15 03:07:10}
\Level{A}
\SubArea{Exponential functions}
\LANGen{
\Question{
Identify a possible graph of the following function. 
\[ 
 f(x) = \left (\frac{2}
{5}\right )^{x+2} - 1
\]}
{}{}{}{}{}{}}
\LANGpl{
\Question{
Wyznacz możliwy wykres podanej funkcji. 
\[ 
 f(x) = \left (\frac{2}
{5}\right )^{x+2} - 1
\]}
{}{}{}{}{}{}}
\LANGcs{
\Question{
Který z následujících grafů je grafem funkce
\(f(x)= \left (\frac{2}
{5}\right )^{x+2} - 1\)?}
{}{}{}{}{}{}}
\LANGsk{
\Question{
Ktorý z následujúcich grafov je grafom funkcie
\(f(x) = \left (\frac{2}
{5}\right )^{x+2} - 1\)?}
{}{}{}{}{}{}}
\LANGes{
\Question{
Identifica la representación gráfica de la siguiente función. 
\[ 
 f(x) = \left (\frac{2}
{5}\right )^{x+2} - 1
\]}
{}{}{}{}{}{}}

\IMGanswer{000036_06-figure2.pdf}
\IMGanswer{000036_06-figure0.pdf}
\IMGanswer{000036_06-figure1.pdf}
\IMGanswer{000036_06-figure3.pdf}
\IMGanswer{000036_06-figure4.pdf}\End

\Data\ID{1003019602}
\Updated{2020-07-15 03:07:12}
\Level{B}
\SubArea{Exponential functions}
\LANGen{
\Question{
Consider values
\[ 4^5;\  0.2^{\frac12};\ \left(\frac54\right)^0;\ \left(\frac13\right)^{-3};\ \left(\frac76\right)^{-3};\ 2.5^{0.6}\text{.} \]
Without using a calculator, determine how many of the values are less than \( 1 \).}
{\( 2 \)}{\( 4 \)}{\( 3 \)}{\( 1 \)}{}{}}
\LANGpl{
\Question{
Rozważ wartości 
\[ 4^5;\  0{,}2^{\frac12};\ \left(\frac54\right)^0;\ \left(\frac13\right)^{-3};\ \left(\frac76\right)^{-3};\ 2{,}5^{0{,}6}\text{.} \]
Nie używając kalkulatora określ ile z podanych wartości jest mniejszych niż \( 1 \).}
{\( 2 \)}{\( 4 \)}{\( 3 \)}{\( 1 \)}{}{}}
\LANGcs{
\Question{
Bez použití kalkulačky určete, kolik z následujících hodnot je menších než jedna.
\[ 4^5;\  0{,}2^{\frac12};\ \left(\frac54\right)^0;\ \left(\frac13\right)^{-3};\ \left(\frac76\right)^{-3};\ 2{,}5^{0{,}6}\]}
{\( 2 \)}{\( 4 \)}{\( 3 \)}{\( 1 \)}{}{}}
\LANGsk{
\Question{
Bez použitia kalkulačky určte, koľko z nasledujúcich hodnôt je menších ako jedna.
\[ 4^5;\  0{,}2^{\frac12};\ \left(\frac54\right)^0;\ \left(\frac13\right)^{-3};\ \left(\frac76\right)^{-3};\ 2{,}5^{0{,}6} \]}
{\( 2 \)}{\( 4 \)}{\( 3 \)}{\( 1 \)}{}{}}
\LANGes{
\Question{
Sin utilizar la calculadora, determina cuántos de estos valores son menores que \( 1 \).
\[ 4^5;\  0.2^{\frac12};\ \left(\frac54\right)^0;\ \left(\frac13\right)^{-3};\ \left(\frac76\right)^{-3};\ 2.5^{0.6}\text{.} \]}
{\( 2 \)}{\( 4 \)}{\( 3 \)}{\( 1 \)}{}{}}
\End

\Data\ID{1003019607}
\Updated{2020-07-15 03:07:12}
\Level{B}
\SubArea{Exponential functions}
\LANGen{
\Question{
Which of the relations between \( m \) and \( n \) ensures \( \left(\frac45\right)^m < \left(\frac45\right)^n \)?}
{\( m > n \)}{\( m \geq n \)}{\( m = n \)}{\( m < n \)}{}{}}
\LANGpl{
\Question{
Która z zależności pomiędzy  \( m \) i \( n \) spełnia \( \left(\frac45\right)^m < \left(\frac45\right)^n \)?}
{\( m > n \)}{\( m \geq n \)}{\( m = n \)}{\( m < n \)}{}{}}
\LANGcs{
\Question{
Rozhodněte, jaký je vztah mezi čísly \( m \) a \( n \), když platí  \( \left(\frac45\right)^m < \left(\frac45\right)^n \).}
{\( m > n \)}{\( m \geq n \)}{\( m = n \)}{\( m < n \)}{}{}}
\LANGsk{
\Question{
Rozhodnite, aký je vzťah medzi číslami \( m \) a \( n \), ak platí  \( \left(\frac45\right)^m < \left(\frac45\right)^n \).}
{\( m > n \)}{\( m \geq n \)}{\( m = n \)}{\( m < n \)}{}{}}
\LANGes{
\Question{
¿Qué relación hay entre \( m \) y \( n \) si \( \left(\frac45\right)^m < \left(\frac45\right)^n \)?}
{\( m > n \)}{\( m \geq n \)}{\( m = n \)}{\( m < n \)}{}{}}
\End

\Data\ID{1003019608}
\Updated{2020-07-15 03:07:12}
\Level{B}
\SubArea{Exponential functions}
\LANGen{
\Question{
Which of the relations between \( m \) and \( n \) ensures \( \left(\frac54\right)^m < \left(\frac54\right)^n \)?}
{\( m < n \)}{\( m \geq n \)}{\( m = n \)}{\( m > n \)}{}{}}
\LANGpl{
\Question{
Która z zależności pomiędzy  \( m \) i \( n \) spełnia \( \left(\frac54\right)^m < \left(\frac54\right)^n \)?}
{\( m < n \)}{\( m \geq n \)}{\( m = n \)}{\( m > n \)}{}{}}
\LANGcs{
\Question{
Rozhodněte, jaký je vztah mezi čísly \( m \) a \( n \), když platí \( \left(\frac54\right)^m < \left(\frac54\right)^n \).}
{\( m < n \)}{\( m \geq n \)}{\( m = n \)}{\( m > n \)}{}{}}
\LANGsk{
\Question{
Rozhodnite, aký je vzťah medzi číslami \( m \) a \( n \), ak platí \( \left(\frac54\right)^m < \left(\frac54\right)^n \).}
{\( m < n \)}{\( m \geq n \)}{\( m = n \)}{\( m > n \)}{}{}}
\LANGes{
\Question{
¿Qué relación hay entre \( m \) y \( n \) si \( \left(\frac54\right)^m < \left(\frac54\right)^n \)?}
{\( m < n \)}{\( m \geq n \)}{\( m = n \)}{\( m > n \)}{}{}}
\End

\Data\ID{1003019609}
\Updated{2021-04-11 09:04:00}
\Level{B}
\SubArea{Exponential functions}
\LANGen{
\Question{
What are the values of the real \( a \), that satisfy inequality \( a^{\frac25} > a^{\frac45}\)?}
{\( 0 < a < 1 \)}{\( a > 1 \)}{\( a < 1 \)}{\( a > 0 \)}{}{}}
\LANGpl{
\Question{
Jakie są wartości rzeczywistego parametru \( a \), które spełniają nierówność \( a^{\frac25} > a^{\frac45}\)?}
{\( 0 < a < 1 \)}{\( a > 1 \)}{\( a < 1 \)}{\( a > 0 \)}{}{}}
\LANGcs{
\Question{
Pro která reálná čísla \( a \)  platí  \( a^{\frac25} > a^{\frac45}\)?}
{\( 0 < a < 1 \)}{\( a > 1 \)}{\( a < 1 \)}{\( a > 0 \)}{}{}}
\LANGsk{
\Question{
Pre ktoré reálne čísla \( a \)  platí  \( a^{\frac25} > a^{\frac45}\)?}
{\( 0 < a < 1 \)}{\( a > 1 \)}{\( a < 1 \)}{\( a > 0 \)}{}{}}
\LANGes{
\Question{
¿Para qué valores reales \( a \) se verifica esta desigualdad \( a^{\frac25} > a^{\frac45}\)?}
{\( 0 < a < 1 \)}{\( a > 1 \)}{\( a < 1 \)}{\( a > 0 \)}{}{}}
\End

\Data\ID{1003019610}
\Updated{2021-04-11 09:04:18}
\Level{B}
\SubArea{Exponential functions}
\LANGen{
\Question{
What are the values of the real \( a \), that satisfy inequality \( a^{-3} > a^{-2}\)?}
{\( 0 < a < 1 \)}{\( a > 1 \)}{\( a < 1 \)}{\( a > 0 \)}{}{}}
\LANGpl{
\Question{
Jakie są wartości rzeczywistego parametru \( a \), które spełniają nierówność  \( a^{-3} > a^{-2}\)?}
{\( 0 < a < 1 \)}{\( a > 1 \)}{\( a < 1 \)}{\( a > 0 \)}{}{}}
\LANGcs{
\Question{
Pro která reálná čísla \( a \) platí \( a^{-3} > a^{-2}\)?}
{\( 0 < a < 1 \)}{\( a > 1 \)}{\( a < 1 \)}{\( a > 0 \)}{}{}}
\LANGsk{
\Question{
Pre ktoré reálne čísla \( a \) platí \( a^{-3} > a^{-2}\)?}
{\( 0 < a < 1 \)}{\( a > 1 \)}{\( a < 1 \)}{\( a > 0 \)}{}{}}
\LANGes{
\Question{
¿Para qué números reales \( a \) se verifica esta desigualdad \( a^{-3} > a^{-2}\)?}
{\( 0 < a < 1 \)}{\( a > 1 \)}{\( a < 1 \)}{\( a > 0 \)}{}{}}
\End

\Data\ID{1003019611}
\Updated{2020-07-15 03:07:12}
\Level{B}
\SubArea{Exponential functions}
\LANGen{
\Question{
Identify which of the following relations is correct.}
{\( \left(\frac15\right)^{-6} > \left(\frac15\right)^{-\frac23} > \left(\frac15\right)^{\frac25} > \left(\frac15\right)^2 > \left(\frac15\right)^7  \)}{\( \left(\frac15\right)^7 > \left(\frac15\right)^2 > \left(\frac15\right)^{\frac25} > \left(\frac15\right)^{-\frac25} > \left(\frac15\right)^{-6}  \)}{\( \left(\frac15\right)^{-\frac23} > \left(\frac15\right)^{-6} > \left(\frac15\right)^{\frac25} > \left(\frac15\right)^2 > \left(\frac15\right)^7  \)}{\( \left(\frac15\right)^{-6} > \left(\frac15\right)^{-\frac23} > \left(\frac15\right)^{\frac25} > \left(\frac15\right)^7 > \left(\frac15\right)^2  \)}{}{}}
\LANGpl{
\Question{
Określ, która z podanych zależności jest poprawna.}
{\( \left(\frac15\right)^{-6} > \left(\frac15\right)^{-\frac23} > \left(\frac15\right)^{\frac25} > \left(\frac15\right)^2 > \left(\frac15\right)^7  \)}{\( \left(\frac15\right)^7 > \left(\frac15\right)^2 > \left(\frac15\right)^{\frac25} > \left(\frac15\right)^{-\frac25} > \left(\frac15\right)^{-6}  \)}{\( \left(\frac15\right)^{-\frac23} > \left(\frac15\right)^{-6} > \left(\frac15\right)^{\frac25} > \left(\frac15\right)^2 > \left(\frac15\right)^7  \)}{\( \left(\frac15\right)^{-6} > \left(\frac15\right)^{-\frac23} > \left(\frac15\right)^{\frac25} > \left(\frac15\right)^7 > \left(\frac15\right)^2  \)}{}{}}
\LANGcs{
\Question{
Rozhodněte, které uspořádání daných hodnot je správné.}
{\( \left(\frac15\right)^{-6} > \left(\frac15\right)^{-\frac23} > \left(\frac15\right)^{\frac25} > \left(\frac15\right)^2 > \left(\frac15\right)^7  \)}{\( \left(\frac15\right)^7 > \left(\frac15\right)^2 > \left(\frac15\right)^{\frac25} > \left(\frac15\right)^{-\frac25} > \left(\frac15\right)^{-6}  \)}{\( \left(\frac15\right)^{-\frac23} > \left(\frac15\right)^{-6} > \left(\frac15\right)^{\frac25} > \left(\frac15\right)^2 > \left(\frac15\right)^7  \)}{\( \left(\frac15\right)^{-6} > \left(\frac15\right)^{-\frac23} > \left(\frac15\right)^{\frac25} > \left(\frac15\right)^7 > \left(\frac15\right)^2  \)}{}{}}
\LANGsk{
\Question{
Rozhodnite, ktoré usporiadanie daných hodnôt je správne.}
{\( \left(\frac15\right)^{-6} > \left(\frac15\right)^{-\frac23} > \left(\frac15\right)^{\frac25} > \left(\frac15\right)^2 > \left(\frac15\right)^7  \)}{\( \left(\frac15\right)^7 > \left(\frac15\right)^2 > \left(\frac15\right)^{\frac25} > \left(\frac15\right)^{-\frac25} > \left(\frac15\right)^{-6}  \)}{\( \left(\frac15\right)^{-\frac23} > \left(\frac15\right)^{-6} > \left(\frac15\right)^{\frac25} > \left(\frac15\right)^2 > \left(\frac15\right)^7  \)}{\( \left(\frac15\right)^{-6} > \left(\frac15\right)^{-\frac23} > \left(\frac15\right)^{\frac25} > \left(\frac15\right)^7 > \left(\frac15\right)^2  \)}{}{}}
\LANGes{
\Question{
Identifica cuál de las siguientes  relaciones es correcta.}
{\( \left(\frac15\right)^{-6} > \left(\frac15\right)^{-\frac23} > \left(\frac15\right)^{\frac25} > \left(\frac15\right)^2 > \left(\frac15\right)^7  \)}{\( \left(\frac15\right)^7 > \left(\frac15\right)^2 > \left(\frac15\right)^{\frac25} > \left(\frac15\right)^{-\frac25} > \left(\frac15\right)^{-6}  \)}{\( \left(\frac15\right)^{-\frac23} > \left(\frac15\right)^{-6} > \left(\frac15\right)^{\frac25} > \left(\frac15\right)^2 > \left(\frac15\right)^7  \)}{\( \left(\frac15\right)^{-6} > \left(\frac15\right)^{-\frac23} > \left(\frac15\right)^{\frac25} > \left(\frac15\right)^7 > \left(\frac15\right)^2  \)}{}{}}
\End

\Data\ID{1003019612}
\Updated{2020-07-15 03:07:12}
\Level{B}
\SubArea{Exponential functions}
\LANGen{
\Question{
Identify which of the following relations is correct.}
{\( 2^7 > 2^2 > 2^{\frac25} > 2^{-\frac25} > 2^{-6}  \)}{\( 2^{-6} > 2^{-\frac23} > 2^{\frac25} > 2^2 > 2^7  \)}{\( 2^{-\frac23} > 2^{-6} > 2^{\frac25} > 2^2 > 2^7  \)}{\( 2^{-6} > 2^{-\frac23} > 2^{\frac25} > 2^7 > 2^2  \)}{}{}}
\LANGpl{
\Question{
Określ, która z poniższych zależności jest poprawna.}
{\( 2^7 > 2^2 > 2^{\frac25} > 2^{-\frac25} > 2^{-6}  \)}{\( 2^{-6} > 2^{-\frac23} > 2^{\frac25} > 2^2 > 2^7  \)}{\( 2^{-\frac23} > 2^{-6} > 2^{\frac25} > 2^2 > 2^7  \)}{\( 2^{-6} > 2^{-\frac23} > 2^{\frac25} > 2^7 > 2^2  \)}{}{}}
\LANGcs{
\Question{
Rozhodněte, které uspořádání daných hodnot je správné.}
{\( 2^7 > 2^2 > 2^{\frac25} > 2^{-\frac25} > 2^{-6}  \)}{\( 2^{-6} > 2^{-\frac23} > 2^{\frac25} > 2^2 > 2^7  \)}{\( 2^{-\frac23} > 2^{-6} > 2^{\frac25} > 2^2 > 2^7  \)}{\( 2^{-6} > 2^{-\frac23} > 2^{\frac25} > 2^7 > 2^2  \)}{}{}}
\LANGsk{
\Question{
Rozhodnite, ktoré usporiadanie daných hodnôt je správne.}
{\( 2^7 > 2^2 > 2^{\frac25} > 2^{-\frac25} > 2^{-6}  \)}{\( 2^{-6} > 2^{-\frac23} > 2^{\frac25} > 2^2 > 2^7  \)}{\( 2^{-\frac23} > 2^{-6} > 2^{\frac25} > 2^2 > 2^7  \)}{\( 2^{-6} > 2^{-\frac23} > 2^{\frac25} > 2^7 > 2^2  \)}{}{}}
\LANGes{
\Question{
Identifica cuál de las siguientes relaciones es correcta.}
{\( 2^7 > 2^2 > 2^{\frac25} > 2^{-\frac25} > 2^{-6}  \)}{\( 2^{-6} > 2^{-\frac23} > 2^{\frac25} > 2^2 > 2^7  \)}{\( 2^{-\frac23} > 2^{-6} > 2^{\frac25} > 2^2 > 2^7  \)}{\( 2^{-6} > 2^{-\frac23} > 2^{\frac25} > 2^7 > 2^2  \)}{}{}}
\End

\Data\ID{1003024901}
\Updated{2021-04-11 09:04:26}
\Level{B}
\SubArea{Exponential functions}
\LANGen{
\Question{
Given the increasing exponential function \(f(x)=a^x\), choose the correct statement about the base \(a\).}
{\(a > 1\)}{\(a=1\)}{\(a < 1\)}{\(  0 < a < 1 \)}{}{}}
\LANGpl{
\Question{
Dana jest funkcja wykładnicza rosnąca  \(f(x)=a^x\). Które z poniższych stwierdzeń dotyczących podstawy \(a\) jest prawdziwe?}
{\(a > 1\)}{\(a=1\)}{\(a < 1\)}{\(  0 < a < 1 \)}{}{}}
\LANGcs{
\Question{
Exponenciální funkce \(f(x)=a^x\) je rostoucí, když pro základ \(a\) platí:}
{\(a > 1\)}{\(a=1\)}{\(a < 1\)}{\(  0 < a < 1 \)}{}{}}
\LANGsk{
\Question{
Exponenciálna funkcia \(f(x)=a^x\) je rastúca, ak pre základ \(a\) platí:}
{\(a > 1\)}{\(a=1\)}{\(a < 1\)}{\(  0 < a < 1 \)}{}{}}
\LANGes{
\Question{
Dada la función exponencial \(f(x)=a^x\), elige para qué valores de \(a\) la función es creciente.}
{\(a > 1\)}{\(a=1\)}{\(a < 1\)}{\(  0 < a < 1 \)}{}{}}
\End

\Data\ID{1003024902}
\Updated{2020-07-15 03:07:12}
\Level{B}
\SubArea{Exponential functions}
\LANGen{
\Question{
Given the function \(f(x)=a\cdot b^x\), where \( a < 0 \) and \( b > 0 \), find the correct statement.}
{Function \( f \) is bounded above.}{Function \( f \) is bounded below.}{Function \( f \) is bounded.}{Function \( f \) is unbounded.}{}{}}
\LANGpl{
\Question{
Dana jest funkcja \(f(x)=a\cdot b^x\), gdzie \( a < 0 \) i \( b > 0 \). Które stwierdzenie jest prawdziwe?}
{Funkcja \( f \) jest ograniczona z góry.}{Funkcja \( f \) jest ograniczona z dołu.}{Funkcja \( f \) jest ograniczona.}{Funkcja \( f \) jest nieograniczona.}{}{}}
\LANGcs{
\Question{
Je dána funkce \(f(x)=a\cdot b^x\), kde \( a < 0 \) a \( b > 0 \). Vyberte pravdivé tvrzení.}
{Funkce \( f \) je omezená shora.}{Funkce \( f \) je omezená zdola.}{Funkce \( f \) je omezená.}{Funkce \( f \) není omezená.}{}{}}
\LANGsk{
\Question{
Daná je funkcia \(f(x)=a\cdot b^x\), kde \( a < 0 \) a \( b > 0 \). Vyberte pravdivé tvrdenie.}
{Funkcia \( f \) je ohraničená zhora.}{Funkcia \( f \) je ohraničená zdola.}{Funkcia \( f \) je ohraničená.}{Funkcia \( f \) je neohraničená.}{}{}}
\LANGes{
\Question{
Dada la función \(f(x)=a\cdot b^x\), donde \( a < 0 \) y \( b > 0 \), halla el enunciado verdadero.}
{La función \( f \) está acotada por arriba.}{La función \( f \) está acotada por abajo.}{La función \( f \) está acotada.}{La función \( f \) no está acotada.}{}{}}
\End

\Data\ID{1003024903}
\Updated{2020-07-15 03:07:12}
\Level{B}
\SubArea{Exponential functions}
\LANGen{
\Question{
Given the function \(f(x)=a\cdot b^x\), where \( a < 0 \) and \( 0 < b < 1 \), find the correct statement.}
{Function \( f \) is increasing.}{Function \( f \) is decreasing.}{Function \( f \) is nonincreasing.}{Function \( f \) is nondecreasing.}{}{}}
\LANGpl{
\Question{
Dana jest funkcja \(f(x)=a\cdot b^x\), gdzie \( a < 0 \) i \( 0 < b < 1 \). Które stwierdzenie jest prawdziwe?}
{Funkcja \( f \) jest rosnąca.}{Funkcja \( f \) jest malejąca.}{Funkcja \( f \) jest nierosnąca.}{Funkcja \( f \) jest niemalejąca.}{}{}}
\LANGcs{
\Question{
Je dána funkce \(f(x)=a\cdot b^x\), kde \( a < 0 \) a \( 0 < b < 1 \). Vyberte pravdivé tvrzení.}
{Funkce \( f \) je rostoucí.}{Funkce \( f \) je klesající.}{Funkce \( f \) je nerostoucí.}{Funkce \( f \) je neklesající.}{}{}}
\LANGsk{
\Question{
Daná je funkcia \(f(x)=a\cdot b^x\), kde \( a < 0 \) a \( 0 < b < 1 \). Vyberte pravdivé tvrdenie.}
{Funkcia \( f \) je rastúca.}{Funkcia \( f \) je klesajúca.}{Funkcia \( f \) je nerastúca.}{Funkcia \( f \) je neklesajúca.}{}{}}
\LANGes{
\Question{
Dada la función \(f(x)=a\cdot b^x\), donde \( a < 0 \) y \( 0 < b < 1 \), halla el enunciado verdadero.}
{La función \( f \) es creciente.}{La función \( f \) es decreciente.}{La función \( f \) es no creciente.}{La función \( f \) es no decreciente.}{}{}}
\End

\Data\ID{1003024904}
\Updated{2020-11-12 01:11:07}
\Level{B}
\SubArea{Exponential functions}
\LANGen{
\Question{
Which of the following functions is decreasing?}
{\( f(x)=\left(\frac15\right)^x \)}{\( f(x)=-5^{-x} \)}{\( f(x)=\left(\frac15\right)^{-x} \)}{\( f(x)=5^x \)}{}{}}
\LANGpl{
\Question{
Która z podanych funkcji jest malejąca?}
{\( f(x)=\left(\frac15\right)^x \)}{\( f(x)=-5^{-x} \)}{\( f(x)=\left(\frac15\right)^{-x} \)}{\( f(x)=5^x \)}{}{}}
\LANGcs{
\Question{
Která z následujících funkcí je klesající?}
{\( f(x)=\left(\frac15\right)^x \)}{\( f(x)=-5^{-x} \)}{\( f(x)=\left(\frac15\right)^{-x} \)}{\( f(x)=5^x \)}{}{}}
\LANGsk{
\Question{
Ktorá z nasledujúcich funkcií je klesajúca?}
{\( f(x)=\left(\frac15\right)^x \)}{\( f(x)=-5^{-x} \)}{\( f(x)=\left(\frac15\right)^{-x} \)}{\( f(x)=5^x \)}{}{}}
\LANGes{
\Question{
¿Cuál de las siguientes funciones es decreciente?}
{\( f(x)=\left(\frac15\right)^x \)}{\( f(x)=-5^{-x} \)}{\( f(x)=\left(\frac15\right)^{-x} \)}{\( f(x)=5^x \)}{}{}}
\End

\Data\ID{1003024905}
\Updated{2020-11-12 01:11:47}
\Level{B}
\SubArea{Exponential functions}
\LANGen{
\Question{
Which of the following functions is bounded above?}
{\( f(x) = -3^{-x} \)}{\( f(x) = \left(\frac13\right)^{-x} \)}{\( f(x) = 3^x \)}{\( f(x) = \left(\frac13\right)^x \)}{}{}}
\LANGpl{
\Question{
Która z podanych funkcji jest ograniczona z góry?}
{\( f(x) = -3^{-x} \)}{\( f(x) = \left(\frac13\right)^{-x} \)}{\( f(x) = 3^x \)}{\( f(x) = \left(\frac13\right)^x \)}{}{}}
\LANGcs{
\Question{
Která z následujících funkcí je omezená shora?}
{\( f(x) = -3^{-x} \)}{\( f(x) = \left(\frac13\right)^{-x} \)}{\( f(x) = 3^x \)}{\( f(x) = \left(\frac13\right)^x \)}{}{}}
\LANGsk{
\Question{
Ktorá z nasledujúcich funkcií je ohraničená zhora?}
{\( f(x) = -3^{-x} \)}{\( f(x) = \left(\frac13\right)^{-x} \)}{\( f(x) = 3^x \)}{\( f(x) = \left(\frac13\right)^x \)}{}{}}
\LANGes{
\Question{
¿Cuál de las siguientes funciones está acotada por arriba?}
{\( f(x) = -3^{-x} \)}{\( f(x) = \left(\frac13\right)^{-x} \)}{\( f(x) = 3^x \)}{\( f(x) = \left(\frac13\right)^x \)}{}{}}
\End

\Data\ID{1003024907}
\Updated{2021-04-11 09:04:55}
\Level{B}
\SubArea{Exponential functions}
\LANGen{
\Question{
Choose the correct list of properties of the function \( f(x)= -\left(\frac12\right)^{-x} \).}
{decreasing, bounded above, asymptote at \( y=0 \)}{decreasing, bounded below, asymptote at \( x=0 \)}{increasing, bounded below, asymptote at \( x=0 \)}{increasing, bounded below, asymptote at \( y=0 \)}{}{}}
\LANGpl{
\Question{
Wybierz odpowiednią listę własności funkcji  \( f(x)= -\left(\frac12\right)^{-x} \).}
{malejąca, ograniczona z góry, asymptota w  \( y=0 \)}{malejąca, ograniczona z  dołu, asymptota w  \( x=0 \)}{rosnąca, ograniczona z dołu, asymptota w  \( x=0 \)}{rosnąca, ograniczona z dołu, asymptota w \( y=0 \)}{}{}}
\LANGcs{
\Question{
Určete, která z následujících možností popisuje vlastnosti funkce \( f(x)= -\left(\frac12\right)^{-x} \).}
{klesající, omezená shora, asymptota \( y=0 \)}{klesající, omezená zdola, asymptota \( x=0 \)}{rostoucí, omezená zdola, asymptota \( x=0 \)}{rostoucí, omezená zdola, asymptota \( y=0 \)}{}{}}
\LANGsk{
\Question{
Ktoré tvrdenie o vlastnostiach funkcie \( f(x)= -\left(\frac12\right)^{-x} \) je pravdivé?}
{klesajúca, ohraničená zhora, asymptota \( y=0 \)}{klesajúca, ohraničená zdola, asymptota \( x=0 \)}{rastúca, ohraničená zdola, asymptota \( x=0 \)}{rastúca, ohraničená zdola, asymptota \( y=0 \)}{}{}}
\LANGes{
\Question{
Elige las propiedades correctas de la función \( f(x)= -\left(\frac12\right)^{-x} \).}
{decreciente, acotada por arriba y con asíntota \( y=0 \)}{decreciente, acotaba por abajo  y con asíntota \( x=0 \)}{creciente, acotada por abajo y conasíntota \( x=0 \)}{creciente, acotada por abajo y con asíntota \( y=0 \)}{}{}}
\End

\Data\ID{1103019601}
\Updated{2021-04-11 09:04:22}
\Level{B}
\SubArea{Exponential functions}
\IMG{1103019601_0.pdf}
\LANGen{
\Question{
Consider values
\[ 0.7^{-0.5};\  \left(\frac58\right)^6;\ \left(\frac32\right)^{-5};\  3.5^0;\  0.4^4;\  5^3\text{.} \]
Determine how many of the values are greater than \( 1 \). Use the graph of an exponential function given below.}
{\( 2 \)}{\( 4 \)}{\( 3 \)}{\( 1 \)}{}{}}
\LANGpl{
\Question{
Dane są wartości 
\[ 0{,}7^{-0{,}5};\  \left(\frac58\right)^6;\ \left(\frac32\right)^{-5};\  3{,}5^0;\  0{,}4^4;\  5^3\text{.} \]
Określ ile z podanych wartości jest większych niż \( 1 \). Skorzystaj z wykresu funkcji wykładniczej, który przedstawiono poniżej.}
{\( 2 \)}{\( 4 \)}{\( 3 \)}{\( 1 \)}{}{}}
\LANGcs{
\Question{
Na obrázku je graf exponenciální funkce. Na základě obrázku rozhodněte, kolik z uvedených  čísel je větších než jedna.
\[ 0{,}7^{-0{,}5};\  \left(\frac58\right)^6;\ \left(\frac32\right)^{-5};\  3{,}5^0;\  0{,}4^4;\  5^3 \]}
{\( 2 \)}{\( 4 \)}{\( 3 \)}{\( 1 \)}{}{}}
\LANGsk{
\Question{
Na obrázku je graf exponenciálnej funkcie. Na základe obrázku rozhodnite, koľko z uvedených  čísel je väčších ako jedna.
\[ 0{,}7^{-0{,}5};\  \left(\frac58\right)^6;\ \left(\frac32\right)^{-5};\  3{,}5^0;\  0{,}4^4;\  5^3\]}
{\( 2 \)}{\( 4 \)}{\( 3 \)}{\( 1 \)}{}{}}
\LANGes{
\Question{
Considera los valores
\[ 0.7^{-0.5};\  \left(\frac58\right)^6;\ \left(\frac32\right)^{-5};\  3.5^0;\  0.4^4;\  5^3\text{.} \]
Determina cuántos valores son mayores que \( 1 \). Usa la gráfica de una función exponencial representada abajo.}
{\( 2 \)}{\( 4 \)}{\( 3 \)}{\( 1 \)}{}{}}
\End

\Data\ID{1103019603}
\Updated{2020-07-15 03:07:12}
\Level{B}
\SubArea{Exponential functions}
\IMG{1103019603_0.pdf}
\LANGen{
\Question{
Identify which of the relations between \( m \) and \(  n \) ensures \( \left(\frac{12}7\right)^m < \left(\frac{12}7\right)^n \). Use the graph of \( f(x) = \left(\frac{12}7\right)^x \) given below.}
{\( m < n \)}{\( m > n \)}{\( m = n \)}{\( m \geq n \)}{}{}}
\LANGpl{
\Question{
Określ, która z zależności pomiędzy  \( m \) i \(  n \) spełnia \( \left(\frac{12}7\right)^m < \left(\frac{12}7\right)^n \). Użyj wykresu funkcji \( f(x) = \left(\frac{12}7\right)^x \) przedstawionego poniżej.}
{\( m < n \)}{\( m > n \)}{\( m = n \)}{\( m \geq n \)}{}{}}
\LANGcs{
\Question{
Na obrázku je graf funkce \( f(x) = \left(\frac{12}7\right)^x \). Jaký je vztah mezi čísly \( m \) a \(  n \), když platí \( \left(\frac{12}7\right)^m < \left(\frac{12}7\right)^n \)?}
{\( m < n \)}{\( m > n \)}{\( m = n \)}{\( m \geq n \)}{}{}}
\LANGsk{
\Question{
Na obrázku je graf funkcie \( f(x) = \left(\frac{12}7\right)^x \). Aký je vzťah medzi číslami \( m \) a \(  n \), ak platí \( \left(\frac{12}7\right)^m < \left(\frac{12}7\right)^n \)?}
{\( m < n \)}{\( m > n \)}{\( m = n \)}{\( m \geq n \)}{}{}}
\LANGes{
\Question{
Identifica qué relación hay entre \( m \) y \(  n \) si \( \left(\frac{12}7\right)^m < \left(\frac{12}7\right)^n \). Usa la gráfica de \( f(x) = \left(\frac{12}7\right)^x \) dada abajo.}
{\( m < n \)}{\( m > n \)}{\( m = n \)}{\( m \geq n \)}{}{}}
\End

\Data\ID{1103019604}
\Updated{2020-11-12 01:11:03}
\Level{B}
\SubArea{Exponential functions}
\IMG{1103019604_0.pdf}
\LANGen{
\Question{
What are the values of the real variable \( x \), if \( 0.7^x \leq 1 \)? Use the graph of \( f(x)=0.7^x \) given below.}
{\( x\in[0;\infty) \)}{\( x\in(-\infty;0) \)}{\( x\in(-\infty;0] \)}{\( x\in(0;\infty) \)}{}{}}
\LANGpl{
\Question{
Jakie są wartości rzeczywistej zmiennej \( x \), jeśli  \( 0{,}7^x \leq 1 \)? Użyj wykresu funkcji \( f(x)=0{,}7^x \) przedstawionego poniżej.}
{\( x\in\langle0;\infty) \)}{\( x\in(-\infty;0) \)}{\( x\in(-\infty;0\rangle \)}{\( x\in(0;\infty) \)}{}{}}
\LANGcs{
\Question{
Na obrázku je graf funkce \( f(x)=0{,}7^x \). Určete všechna reálná čísla  \( x \), pro která platí \( 0{,}7^x \leq 1 \).}
{\( x\in\langle0;\infty) \)}{\( x\in(-\infty;0) \)}{\( x\in(-\infty;0\rangle \)}{\( x\in(0;\infty) \)}{}{}}
\LANGsk{
\Question{
Na obrázku je graf funkcie \( f(x)=0{,}7^x \). Určte všetky reálne čísla  \( x \), pre ktoré platí \( 0{,}7^x \leq 1 \).}
{\( x\in\langle0;\infty) \)}{\( x\in(-\infty;0) \)}{\( x\in(-\infty;0\rangle \)}{\( x\in(0;\infty) \)}{}{}}
\LANGes{
\Question{
En la imagen está la gráfica de la función \( f(x)=0.7^x \). Determina todos los valores del número real \( x \) si \( 0.7^x \leq 1 \)?}
{\( x\in[0;\infty) \)}{\( x\in(-\infty;0) \)}{\( x\in(-\infty;0] \)}{\( x\in(0;\infty) \)}{}{}}
\End

\Data\ID{1103019605}
\Updated{2020-11-12 01:11:16}
\Level{B}
\SubArea{Exponential functions}
\IMG{1103019605_0.pdf}
\LANGen{
\Question{
What are the values of the real variable \( x \), if \( 5^x \leq 1 \)? Use the graph of \( f(x)=5^x \) given below.}
{\( x\in (-\infty;0] \)}{\( x\in [0;\infty) \)}{\( x\in (-\infty;0) \)}{\( x\in (0;\infty) \)}{}{}}
\LANGpl{
\Question{
Jakie są wartości rzeczywistej  zmiennej \( x \), jeśli \( 5^x \leq 1 \)? Wykorzystaj wykres funkcji  \( f(x)=5^x \) przedstawiony poniżej.}
{\( x\in (-\infty;0\rangle \)}{\( x\in \langle0;\infty) \)}{\( x\in (-\infty;0) \)}{\( x\in (0;\infty) \)}{}{}}
\LANGcs{
\Question{
Na obrázku je graf funkce \( f(x)=5^x \). Určete všechna reálná čísla \( x \), pro která platí  \( 5^x \leq 1 \).}
{\( x\in (-\infty;0\rangle \)}{\( x\in \langle0;\infty) \)}{\( x\in (-\infty;0) \)}{\( x\in (0;\infty) \)}{}{}}
\LANGsk{
\Question{
Na obrázku je graf funkcie \( f(x)=5^x \). Určte všetky reálne čísla \( x \), pre ktoré platí  \( 5^x \leq 1 \).}
{\( x\in (-\infty;0\rangle \)}{\( x\in \langle0;\infty) \)}{\( x\in (-\infty;0) \)}{\( x\in (0;\infty) \)}{}{}}
\LANGes{
\Question{
En la imagen está la gráfica de la función \( f(x)=5^x \). Determina todos los valores del número real \( x \) si \( 5^x \leq 1 \)?}
{\( x\in (-\infty;0] \)}{\( x\in [0;\infty) \)}{\( x\in (-\infty;0) \)}{\( x\in (0;\infty) \)}{}{}}
\End

\Data\ID{1103019606}
\Updated{2020-07-15 03:07:12}
\Level{B}
\SubArea{Exponential functions}
\IMG{1103019606_0.pdf}
\LANGen{
\Question{
Identify which of the following relations is correct.  Use the graph of \( f(x)=\left(\frac12\right)^x \) given below.}
{\( \left(\frac12\right)^{-6} >\left(\frac12\right)^{-\frac23} >\left(\frac12\right)^{\frac25} >\left(\frac12\right)^2 >\left(\frac12\right)^7 \)}{\( \left(\frac12\right)^7 >\left(\frac12\right)^2 >\left(\frac12\right)^{\frac25} >\left(\frac12\right)^{-\frac25}>\left(\frac12\right)^{-6} \)}{\( \left(\frac12\right)^{-\frac23} >\left(\frac12\right)^{-6} >\left(\frac12\right)^{\frac25} >\left(\frac12\right)^2 >\left(\frac12\right)^7 \)}{\( \left(\frac12\right)^{-6} >\left(\frac12\right)^{-\frac23} >\left(\frac12\right)^{\frac25} >\left(\frac12\right)^7 >\left(\frac12\right)^2 \)}{}{}}
\LANGpl{
\Question{
Która z podanych zależności jest prawdziwa?  Użyj wykresu funkcji \( f(x)=\left(\frac12\right)^x \) przedstawionego poniżej.}
{\( \left(\frac12\right)^{-6} >\left(\frac12\right)^{-\frac23} >\left(\frac12\right)^{\frac25} >\left(\frac12\right)^2 >\left(\frac12\right)^7 \)}{\( \left(\frac12\right)^7 >\left(\frac12\right)^2 >\left(\frac12\right)^{\frac25} >\left(\frac12\right)^{-\frac25}>\left(\frac12\right)^{-6} \)}{\( \left(\frac12\right)^{-\frac23} >\left(\frac12\right)^{-6} >\left(\frac12\right)^{\frac25} >\left(\frac12\right)^2 >\left(\frac12\right)^7 \)}{\( \left(\frac12\right)^{-6} >\left(\frac12\right)^{-\frac23} >\left(\frac12\right)^{\frac25} >\left(\frac12\right)^7 >\left(\frac12\right)^2 \)}{}{}}
\LANGcs{
\Question{
Na obrázku je graf funkce \( f(x)=\left(\frac12\right)^x \). Rozhodněte, které uspořádání hodnot je správné.}
{\( \left(\frac12\right)^{-6} >\left(\frac12\right)^{-\frac23} >\left(\frac12\right)^{\frac25} >\left(\frac12\right)^2 >\left(\frac12\right)^7 \)}{\( \left(\frac12\right)^7 >\left(\frac12\right)^2 >\left(\frac12\right)^{\frac25} >\left(\frac12\right)^{-\frac25}>\left(\frac12\right)^{-6} \)}{\( \left(\frac12\right)^{-\frac23} >\left(\frac12\right)^{-6} >\left(\frac12\right)^{\frac25} >\left(\frac12\right)^2 >\left(\frac12\right)^7 \)}{\( \left(\frac12\right)^{-6} >\left(\frac12\right)^{-\frac23} >\left(\frac12\right)^{\frac25} >\left(\frac12\right)^7 >\left(\frac12\right)^2 \)}{}{}}
\LANGsk{
\Question{
Na obrázku je graf funkcie \( f(x)=\left(\frac12\right)^x \). Rozhodnite, ktoré usporiadanie hodnôt je správne.}
{\( \left(\frac12\right)^{-6} >\left(\frac12\right)^{-\frac23} >\left(\frac12\right)^{\frac25} >\left(\frac12\right)^2 >\left(\frac12\right)^7 \)}{\( \left(\frac12\right)^7 >\left(\frac12\right)^2 >\left(\frac12\right)^{\frac25} >\left(\frac12\right)^{-\frac25}>\left(\frac12\right)^{-6} \)}{\( \left(\frac12\right)^{-\frac23} >\left(\frac12\right)^{-6} >\left(\frac12\right)^{\frac25} >\left(\frac12\right)^2 >\left(\frac12\right)^7 \)}{\( \left(\frac12\right)^{-6} >\left(\frac12\right)^{-\frac23} >\left(\frac12\right)^{\frac25} >\left(\frac12\right)^7 >\left(\frac12\right)^2 \)}{}{}}
\LANGes{
\Question{
En la imagen está la gráfica de la función \( f(x)=\left(\frac12\right)^x \). Identifica cuál de las siguientes relaciones es correcta.}
{\( \left(\frac12\right)^{-6} >\left(\frac12\right)^{-\frac23} >\left(\frac12\right)^{\frac25} >\left(\frac12\right)^2 >\left(\frac12\right)^7 \)}{\( \left(\frac12\right)^7 >\left(\frac12\right)^2 >\left(\frac12\right)^{\frac25} >\left(\frac12\right)^{-\frac25}>\left(\frac12\right)^{-6} \)}{\( \left(\frac12\right)^{-\frac23} >\left(\frac12\right)^{-6} >\left(\frac12\right)^{\frac25} >\left(\frac12\right)^2 >\left(\frac12\right)^7 \)}{\( \left(\frac12\right)^{-6} >\left(\frac12\right)^{-\frac23} >\left(\frac12\right)^{\frac25} >\left(\frac12\right)^7 >\left(\frac12\right)^2 \)}{}{}}
\End

\Data\ID{1103024906}
\Updated{2021-04-11 09:04:09}
\Level{B}
\SubArea{Exponential functions}
\IMG{1103024906_0.pdf}
\LANGen{
\Question{
The graph of function \(-3^{-x}\) is below. Choose the correct list of properties of \( f\).}
{increasing, bounded above, asymptote at  \( y=0 \)}{decreasing, bounded above, asymptote at  \( y=0 \)}{decreasing, bounded below, asymptote at  \( x=0 \)}{increasing, bounded below, asymptote at  \( x=0 \)}{}{}}
\LANGpl{
\Question{
Wykres funkcji \(-3^{-x}\) przedstawiono poniżej. Które z niżej podanych są własnościami funkcji \( f\)?}
{rosnąca, ograniczona z góry, asymptota w  \( y=0 \)}{malejąca, ograniczona z góry, asymptota w \( y=0 \)}{malejąca, ograniczona  z dołu, asymptota w \( x=0 \)}{rosnąca, ograniczona z dołu, asymptota w   \( x=0 \)}{}{}}
\LANGcs{
\Question{
Na obrázku je graf funkce \(-3^{-x}\). Rozhodněte, která z následujících možností popisuje vlastnosti funkce f.}
{rostoucí, omezená shora, asymptota \( y=0 \)}{klesající, omezená shora, asymptota \( y=0 \)}{klesající, omezená zdola, asymptota \( x=0 \)}{rostoucí, omezená zdola, asymptota \( x=0 \)}{}{}}
\LANGsk{
\Question{
Na obrázku je graf funkcie \(-3^{-x}\). Vyberte pravdivé tvrdenie o vlastnostiach funkcie \( f\).}
{rastúca, ohraničená zhora, asymptota \( y=0 \)}{klesajúca, ohraničená zhora, asymptota \( y=0 \)}{klesajúca, ohraničená zdola, asymptota \( x=0 \)}{rastúca, ohraničená zdola, asymptota \( x=0 \)}{}{}}
\LANGes{
\Question{
En la imagen está la función \(-3^{-x}\). Elige cuál de las propiedades de la función \( f\) es correcta.}
{creciente, acotada por arriba, asíntota en \( y=0 \)}{decreciente, acotada por arriba, asíntota en \( y=0 \)}{decreciente, acotada por abajo, asíntota en \( x=0 \)}{creciente, acotada por abajo, asíntota en \( x=0 \)}{}{}}
\End

\Data\ID{1103024908}
\Updated{2021-04-11 09:04:53}
\Level{B}
\SubArea{Exponential functions}
\LANGen{
\Question{
The graphs of functions \(f(x)=a\cdot 2^{bx}+2\), where \(a\in\{-1,1\}\), \(b\in\{-1,1\}\), are below. Identify which of the graphs represents the function that is increasing, bounded below, and has an asymptote at \(y=2\).}
{}{}{}{}{}{}}
\LANGpl{
\Question{
Wykresy funkcji \(f(x)=a\cdot 2^{bx}+2\), gdzie \(a\in\{-1,1\}\), \(b\in\{-1,1\}\), przedstawiono poniżej. Który z poniższych wykresów przedstawia funkcję, która jest rosnąca, ograniczona z dołu, i ma asymptotę w  \(y=2\).}
{}{}{}{}{}{}}
\LANGcs{
\Question{
Na následujících obrázcích jsou grafy funkcí \(f(x)=a\cdot 2^{bx}+2\), kde \(a\in\{-1,1\}\), \(b\in\{-1,1\}\). Vyberte graf funkce, která je rostoucí, zdola omezená a má asymptotu \(y=2\).}
{}{}{}{}{}{}}
\LANGsk{
\Question{
Na nasledujúcich obrázkoch sú grafy funkcií \(f(x)=a\cdot 2^{bx}+2\), kde \(a\in\{-1,1\}\), \(b\in\{-1,1\}\). Vyberte graf funkcie, ktorá je rastúca, zdola ohraničená a má asymptotu \(y=2\).}
{}{}{}{}{}{}}
\LANGes{
\Question{
En las imágenes están las funciónes  \(f(x)=a\cdot 2^{bx}+2\), dónde \(a\in\{-1,1\}\), \(b\in\{-1,1\}\). Identifica cuál de las gráficas representa la función que es creciente, acotada inferiormente y tiene una asíntota en \(y=2\).}
{}{}{}{}{}{}}

\IMGanswer{1103024908a.pdf}
\IMGanswer{1103024908b.pdf}
\IMGanswer{1103024908c_0.pdf}
\IMGanswer{1103024908d_0.pdf}\End

\Data\ID{9000003602}
\Updated{2020-07-15 03:07:34}
\Level{B}
\SubArea{Exponential functions}
\LANGen{
\Question{
Which values of the real parameter
\(p\) ensure that
the function \(f(x) = \left (\frac{p+1}
{p-3}\right )^{x}\) 
is an increasing function?}
{\(p\in (3;\infty )\)}{\(p\in \mathbb{R}\)}{\(p\in \mathbb{R}\setminus \{3\}\)}{\(p\in (-\infty ;-1)\cup (3;\infty )\)}{}{}}
\LANGpl{
\Question{
Które wartości rzeczywistego parametru
\(p\), dla których funkcja  \(f(x) = \left (\frac{p+1}
{p-3}\right )^{x}\) 
jest funkcją rosnącą?}
{\(p\in (3;\infty )\)}{\(p\in \mathbb{R}\)}{\(p\in \mathbb{R}\setminus \{3\}\)}{\(p\in (-\infty ;-1)\cup (3;\infty )\)}{}{}}
\LANGcs{
\Question{
Určete všechny hodnoty reálného parametru
\(p\) tak, aby
funkce \(f(x) = \left (\frac{p+1}
{p-3}\right )^{x}\) 
byla rostoucí.}
{\(p\in (3;\infty )\)}{\(p\in \mathbb{R}\)}{\(p\in \mathbb{R}\setminus \{3\}\)}{\(p\in (-\infty ;-1)\cup (3;\infty )\)}{}{}}
\LANGsk{
\Question{
Rozhodnite, pre ktoré hodnoty reálneho parametra
\(p\) je
funkcia \(f(x) = \left (\frac{p+1}
{p-3}\right )^{x}\) 
rastúca.}
{\(p\in (3;\infty )\)}{\(p\in \mathbb{R}\)}{\(p\in \mathbb{R}\setminus \{3\}\)}{\(p\in (-\infty ;-1)\cup (3;\infty )\)}{}{}}
\LANGes{
\Question{
Determina todos los valores del parámetro real 
\(p\) para que la función
\(f(x) = \left (\frac{p+1}
{p-3}\right )^{x}\) 
sea creciente.}
{\(p\in (3;\infty )\)}{\(p\in \mathbb{R}\)}{\(p\in \mathbb{R}\setminus \{3\}\)}{\(p\in (-\infty ;-1)\cup (3;\infty )\)}{}{}}
\End

\Data\ID{9000003603}
\Updated{2021-04-11 09:04:04}
\Level{B}
\SubArea{Exponential functions}
\LANGen{
\Question{
Using properties of exponential function convert the
following inequality to an explicit inequality for the parameter
\(a\).
\[ 
 \left (\sqrt{3} -\sqrt{2}\right )^{2a+1} > \left (\sqrt{3} -\sqrt{2}\right )^{4-a}
\]}
{\(a < 1\)}{\(a > 0\)}{\(0 < a < 1\)}{\(a > 1\)}{}{}}
\LANGpl{
\Question{
Używając własności funkcji wykładniczej przekształć podaną nierówność na nierówność jednoznaczną dla parametru 
\(a\).
\[ 
 \left (\sqrt{3} -\sqrt{2}\right )^{2a+1} > \left (\sqrt{3} -\sqrt{2}\right )^{4-a}
\]}
{\(a < 1\)}{\(a > 0\)}{\(0 < a < 1\)}{\(a > 1\)}{}{}}
\LANGcs{
\Question{
Užitím vlastností exponenciální funkce doplňte tvrzení: Pro reálné číslo \(a\) 
platí vztah
\[ 
 \left (\sqrt{3} -\sqrt{2}\right )^{2a+1} > \left (\sqrt{3} -\sqrt{2}\right )^{4-a}
\]
 právě tehdy, když}
{\(a < 1\).}{\(a > 0\).}{\(0 < a < 1\).}{\(a > 1\).}{}{}}
\LANGsk{
\Question{
Využitím vlastností exponenciálnej funkcie doplňte tvrdenie: Pre reálne číslo \(a\) 
platí vzťah
\[ 
 \left (\sqrt{3} -\sqrt{2}\right )^{2a+1} > \left (\sqrt{3} -\sqrt{2}\right )^{4-a}
\]
 práve vtedy, keď}
{\(a < 1\).}{\(a > 0\).}{\(0 < a < 1\).}{\(a > 1\).}{}{}}
\LANGes{
\Question{
Usando las propiedades de la función exponencial, convierte la siguiente desigualdad en desigualdad estricta para el parametro 
\(a\).
\[ 
 \left (\sqrt{3} -\sqrt{2}\right )^{2a+1} > \left (\sqrt{3} -\sqrt{2}\right )^{4-a}
\]}
{\(a < 1\)}{\(a > 0\)}{\(0 < a < 1\)}{\(a > 1\)}{}{}}
\End

\Data\ID{9000003704}
\Updated{2020-07-15 03:07:34}
\Level{B}
\SubArea{Exponential functions}
\IMG{000037_04-figure0.pdf}
\LANGen{
\Question{
The function \(g(x) = 3 - 3^{x}\) 
is graphed in the picture. In the following list identify one statement which is not
true.}
{The range of the function \(g\) 
is \((-\infty ;3] \).}{The function \(g\) 
is neither odd nor even.}{The function \(g\) 
is decreasing on the domain.}{The domain of the function \(g\) 
is \((-\infty ;\infty )\).}{The function \(g\) 
is not bounded. It is bounded above.}{All the values of the function \(g\) 
are smaller than \(3\).}}
\LANGpl{
\Question{
Funkcja  \(g(x) = 3 - 3^{x}\) 
jest przedstawiona poniżej za pomocą wykresu. Określ, które z poniższych stwierdzeń nie jest prawdziwe.}
{Zakres funkcji  \(g\) 
wynosi  \((-\infty ;3] \).}{Funkcja   \(g\) 
nie jest ani nieparzysta ani parzysta.}{Funkcja \(g\) 
jest malejąca w dziedzinie.}{Dziedziną funkcji \(g\) 
jest  \((-\infty ;\infty )\).}{Funkcja  \(g\) 
jest ograniczona z góry.}{Wszystkie wartości funkcji \(g\) 
są mniejsze od \(3\).}}
\LANGcs{
\Question{
Je dána funkce \(g(x) = 3 - 3^{x}\) 
(viz obrázek). Z následujících tvrzení vyberte to, které není
pravdivé.}
{Obor hodnot funkce je interval \((-\infty ;3\rangle \).}{Funkce není sudá ani lichá.}{Funkce \(g\) je na svém definičním oboru klesající.}{Definiční obor funkce \(g\) je \((-\infty ;\infty )\).}{Funkce je shora omezená, ale není omezená.}{Funkce má všechny funkční hodnoty menší než
\(3\).}}
\LANGsk{
\Question{
Je daná funkcia \(g(x) = 3 - 3^{x}\) 
na obrázku. Ktoré z následujúcich tvrdení nie je pravdivé?}
{Obor hodnôt funkcie je interval \((-\infty ;3\rangle \).}{Funkcia nie je párna nie je nepárna.}{Funkcia \(g\)  je na svojom definičnom obore klesajúca.}{Definičný obor funkcie \(g \) je \((-\infty ;\infty )\).}{Funkcia je zhora ohraničená, ale nie je ohraničená.}{Funkcia má všetky funkčné hodnoty menšie než
\(3\).}}
\LANGes{
\Question{
Dada la función \(g(x) = 3 - 3^{x}\) (mira la imagen). Identifica cuál de los enunciados no es verdadero.}
{El rango de la función \(g\) 
es \((-\infty ;3] \).}{La función \(g\) no es ni par, ni impar.}{La función \(g\) 
es decreciente en su dominio.}{El dominio de la función \(g\) 
es \((-\infty ;\infty )\).}{La función \(g\) 
no está acotada. Está acotada por arriba.}{Todos los valores de la función \(g\) 
son menores que \(3\).}}
\End

\Data\ID{1003025503}
\Updated{2021-04-11 09:04:19}
\Level{C}
\SubArea{Exponential functions}
\LANGen{
\Question{
Find the domain of the function \( f(x)=\sqrt{3^x}\).}
{\( (-\infty;\infty) \)}{\( (0;\infty) \)}{\( [0;\infty) \)}{\( (-\infty;0) \)}{}{}}
\LANGpl{
\Question{
Wyznacz dziedzinę funkcji \( f(x)=\sqrt{3^x}\).}
{\( (-\infty;\infty) \)}{\( (0;\infty) \)}{\( \langle0;\infty) \)}{\( (-\infty;0) \)}{}{}}
\LANGcs{
\Question{
Určete definiční obor funkce \( f(x)=\sqrt{3^x}\).}
{\( (-\infty;\infty) \)}{\( (0;\infty) \)}{\( \langle0;\infty) \)}{\( (-\infty;0) \)}{}{}}
\LANGsk{
\Question{
Určte definičný obor funkcie \( f(x)=\sqrt{3^x}\).}
{\( (-\infty;\infty) \)}{\( (0;\infty) \)}{\( \langle0;\infty) \)}{\( (-\infty;0) \)}{}{}}
\LANGes{
\Question{
Halla el dominio de la función \( f(x)=\sqrt{3^x}\).}
{\( (-\infty;\infty) \)}{\( (0;\infty) \)}{\( [0;\infty) \)}{\( (-\infty;0) \)}{}{}}
\End

\Data\ID{1003025504}
\Updated{2021-04-11 09:04:32}
\Level{C}
\SubArea{Exponential functions}
\LANGen{
\Question{
Find the range of the function \( f(x)=\Bigl|\left(\frac13\right)^x-1\Bigr|\).}
{\( [ 0;\infty) \)}{\( ( 0;\infty) \)}{\( (-1;\infty) \)}{\( [ -1;\infty) \)}{}{}}
\LANGpl{
\Question{
Określ zakres funkcji \( f(x)=\Bigl|\left(\frac13\right)^x-1\Bigr|\).}
{\( \langle 0;\infty) \)}{\( ( 0;\infty) \)}{\( (-1;\infty) \)}{\( \langle -1;\infty) \)}{}{}}
\LANGcs{
\Question{
Určete obor hodnot funkce \( f(x)=\Bigl|\left(\frac13\right)^x-1\Bigr|\).}
{\( \langle 0;\infty) \)}{\( ( 0;\infty) \)}{\( (-1;\infty) \)}{\( \langle -1;\infty) \)}{}{}}
\LANGsk{
\Question{
Určte obor hodnôt funkcie \( f(x)=\Bigl|\left(\frac13\right)^x-1\Bigr|\).}
{\( \langle 0;\infty) \)}{\( ( 0;\infty) \)}{\( (-1;\infty) \)}{\( \langle -1;\infty) \)}{}{}}
\LANGes{
\Question{
Halla el rango de la función \( f(x)=\Bigl|\left(\frac13\right)^x-1\Bigr|\).}
{\( [ 0;\infty) \)}{\( ( 0;\infty) \)}{\( (-1;\infty) \)}{\( [ -1;\infty) \)}{}{}}
\End

\Data\ID{1003101102}
\Updated{2020-07-15 03:07:12}
\Level{C}
\SubArea{Exponential functions}
\LANGen{
\Question{
Find the false statement about the function  \( f(x)=\left(\frac12\right)^{|x|} \).}
{The function \( f \) has the minimum at \( x=0 \).}{The function \( f \) has the maximum at \( x=0 \).}{The function \( f \) is bounded.}{The function \( f \) is even.}{}{}}
\LANGpl{
\Question{
Wybierz fałszywe zdanie dotyczące funkcji  \( f(x)=\left(\frac12\right)^{|x|} \).}
{Funkcja \( f \) osiąga minimum w \( x=0 \).}{Funkcja \( f \) osiąga maksimum w \( x=0 \).}{Funkcja \( f \) jest ograniczona.}{Funkcja \( f \) jest parzysta.}{}{}}
\LANGcs{
\Question{
Funkce \( f \) je dána předpisem \( f(x)=\left(\frac12\right)^{|x|} \). Vyberte nepravdivý výrok.}
{Funkce \( f \) má minimum v bodě \( x=0 \).}{Funkce \( f \) má maximum v bodě \( x=0 \).}{Funkce \( f \) je omezená.}{Funkce \( f \) je sudá.}{}{}}
\LANGsk{
\Question{
Daná je funkcia \( f(x)=\left(\frac12\right)^{|x|} \). Vyberte nesprávne tvrdenie o funkcii \( f \).}
{Funkcia \( f \) má minimum v bode \( x=0 \).}{Funkcia \( f \) má maximum v bode \( x=0 \).}{Funkcia \( f \) je ohraničená.}{Funkcia \( f \) je párna.}{}{}}
\LANGes{
\Question{
Halla el enunciado falso sobre la función \( f(x)=\left(\frac12\right)^{|x|} \).}
{La función \( f \) tiene el mínimo en el punto \( x=0 \).}{La función \( f \) tiene el máximo en el punto \( x=0 \).}{La función \( f \) está acotada.}{La función \( f \) es par.}{}{}}
\End

\Data\ID{1103025501}
\Updated{2020-07-15 03:07:10}
\Level{C}
\SubArea{Exponential functions}
\LANGen{
\Question{
Identify a possible graph of the function \( f(x)= 2^{|x|}-1\).}
{}{}{}{}{}{}}
\LANGpl{
\Question{
Który z poniższych wykresów może być wykresem funkcji \( f(x)= 2^{|x|}-1\)?}
{}{}{}{}{}{}}
\LANGcs{
\Question{
Který z následujících grafů je grafem funkce \( f(x)= 2^{|x|}-1\)?}
{}{}{}{}{}{}}
\LANGsk{
\Question{
Ktorý z nasledujúcich grafov je grafom funkcie \( f(x)= 2^{|x|}-1\)?}
{}{}{}{}{}{}}
\LANGes{
\Question{
Identifica la representación gráfica de la función \( f(x)= 2^{|x|}-1\).}
{}{}{}{}{}{}}

\IMGanswer{1103025501a_0.pdf}
\IMGanswer{1103025501b_0.pdf}
\IMGanswer{1103025501c_0.pdf}
\IMGanswer{1103025501d_0.pdf}\End

\Data\ID{1103025502}
\Updated{2020-07-15 03:07:10}
\Level{C}
\SubArea{Exponential functions}
\LANGen{
\Question{
Identify a possible graph of the function \( f(x)=\Bigl|\left(\frac12\right)^x-1\Bigr|\).}
{}{}{}{}{}{}}
\LANGpl{
\Question{
Który z podanych wykresów może być wykresem funkcji \( f(x)=\Bigl|\left(\frac12\right)^x-1\Bigr|\)?}
{}{}{}{}{}{}}
\LANGcs{
\Question{
Který z následujících grafů je grafem funkce \( f(x)=\Bigl|\left(\frac12\right)^x-1\Bigr|\)?}
{}{}{}{}{}{}}
\LANGsk{
\Question{
Ktorý z nasledujúcich grafov je grafom funkcie \( f(x)=\Bigl|\left(\frac12\right)^x-1\Bigr|\)?}
{}{}{}{}{}{}}
\LANGes{
\Question{
Identifica la representación gráfica de la función \( f(x)=\Bigl|\left(\frac12\right)^x-1\Bigr|\).}
{}{}{}{}{}{}}

\IMGanswer{1103025502a_0.pdf}
\IMGanswer{1103025502b_0.pdf}
\IMGanswer{1103025502c_0.pdf}
\IMGanswer{1103025502d_0.pdf}\End

\Data\ID{1103082501}
\Updated{2021-04-11 09:04:58}
\Level{C}
\SubArea{Exponential functions}
\IMG{1103082501.pdf}
\LANGen{
\Question{
Mary purchased a new car at the purchase price \( 12\,000 \) EUR. Depreciation rate of the car is \( 12\% \) per year. Suzy purchased a new car at the same time as Mary and the purchase price was \( 14\,500\) EUR. The depreciation rate of Suzy's car is \( 15\% \) per year. Let \( p \) be the price of the car in thousands euros and \( t \) be the age of the car in years. Which of the following legends is correct for the next graph?}
{Suzy's car, \( p=14.5\cdot(0.85)^t \)}{Mary's car, \( p=12.0\cdot(0.88)^t \)}{Suzy's car, \( p=14.5\cdot(1.15)^t \)}{Mary's car, \( p=12.0\cdot(1.12)^t \)}{}{}}
\LANGpl{
\Question{
Marysia zakupiła samochód za \(  12\,000 \) EUR. Wskaźnik spadku wartości samochodu w skali rocznej wynosi  \( 12\% \). Zuza kupiła nowy samochód w tym samym czasie co Marysia i zapłaciła za niego \( 14\,500\) EUR. Wskaźnik spadku wartości samochodu Zuzy w skali rocznej wynosi\( 15\% \). Niech \( p \) ceną samochodu wyrażoną w  tysiącach euro, a  \( t \) wiekiem samochodu wyrażonym w latach. Która z poniższych legend do podanego wykresu  jest poprawna?}
{Samochód Zuzy, \( p=14{,}5\cdot(0{,}85)^t \)}{Samochód Marysi, \( p=12{,}0\cdot(0{,}88)^t \)}{Samochód Zuzy, \( p=14{,}5\cdot(1{,}15)^t \)}{Samochód Marysi, \( p=12{,}0\cdot(1{,}12)^t \)}{}{}}
\LANGcs{
\Question{
Rebeka si koupila nové auto za cenu \(  12\,000 \) EUR.  Deprese (pokles) hodnoty auta je \( 12\% \) ročně.  Zuzka si také ve stejné době jako Rebeka koupila nové auto za cenu \(14\,500\) EUR. Deprese (pokles) hodnoty Zuzčina auta je \( 15\% \) ročně. Nechť \( p \) je cena auta v tisících eur a \( t \) je věk auta v rocích. Která z následujících možností správně popisuje daný graf?}
{Zuzčino auto, \( p=14{,}5\cdot(0{,}85)^t \)}{Rebečino auto, \( p=12{,}0\cdot(0{,}88)^t \)}{Zuzčino auto, \( p=14{,}5\cdot(1{,}15)^t \)}{Rebečino auto, \( p=12{,}0\cdot(1{,}12)^t \)}{}{}}
\LANGsk{
\Question{
Rebeka si kúpila nové auto za cenu \( 12\,000 \) EUR.  Depreciácia (pokles) hodnoty auta je \( 12\% \) ročne.  Zuzka si tiež v tom istom čase ako Rebeka kúpila nové auto za cenu \(14\,500\) EUR. Depreciácia (pokles) hodnoty Zuzkinho auta je \( 15\% \) ročne. Nech \( p \) je cena auta v tisícoch eur a \( t \) je vek auta v rokoch. Ktorá z nasledujúcich možností správne popisuje daný graf?}
{Zuzkine auto, \( p=14{,}5\cdot(0{,}85)^t \)}{Rebekine auto, \( p=12{,}0\cdot(0{,}88)^t \)}{Zuzkine auto, \( p=14{,}5\cdot(1{,}15)^t \)}{Rebekine auto, \( p=12{,}0\cdot(1{,}12)^t \)}{}{}}
\LANGes{
\Question{
María compró un coche nuevo por \( 12\,000 \) EUR. La disminución de valor del coche es del \( 12\% \) al año. Susana compró un coche nuevo el mismo año que María y pagó \( 14\,500\) EUR. La disminución del valor del coche de Susana es del  \( 15\% \) al año. Sea \( p \) el precio de coche en miles de euros y \( t \) la el tiempo del coche en años. ¿Cuál de las siguientes posibilidades es correcta según la gráfica representada?}
{El coche de Susana, \( p=14.5\cdot(0.85)^t \)}{El coche de María, \( p=12.0\cdot(0.88)^t \)}{El coche de Susana, \( p=14.5\cdot(1.15)^t \)}{El coche de María, \( p=12.0\cdot(1.12)^t \)}{}{}}
\End

\Data\ID{1103082502}
\Updated{2021-04-11 09:04:16}
\Level{C}
\SubArea{Exponential functions}
\IMG{1103082502.pdf}
\LANGen{
\Question{
Mary purchased a new car. The purchase price was \(12\,000\) euros and the value of the car depreciates by \( 12\% \) per year. Suzy purchased a new car at the same time as Mary and the purchase price was \( 14\,500 \) euros. The depreciation rate for Suzy's car is \( 15\% \).  Let \( p \) be the price of the car in thousands euros and \( t \) be the age of the car in years. In the diagram below, determine the color of the graph that shows relationship between the price of Mary's car and its age. (Note: The answers include not only the color of the graph but also the relevant function equation.)}
{green, \( p=12.0\cdot(0.88)^t \)}{yellow, \( p=14.5\cdot(1.15)^t \)}{orange, \( p=12.0\cdot(1.12)^t \)}{blue, \( p=12.0\cdot(0.88)^t \)}{}{}}
\LANGpl{
\Question{
Marysia kupiła nowy samochód. Cena zakupu wyniosła \(12\,000\) euro, a jego wartość spada o \( 12\% \) rocznie. Zuza kupiła nowy samochód w tym samym czasie co Marysia i zapłaciła za niego \( 14\,500 \) euro. Wskaźnik spadku wartości samochodu Suzy wynosi \( 15\% \).  Niech \( p \) ceną samochodu wyrażona w tysiącach euro, a  \( t \) wiekiem samochodu wyrażonym w latach. Na poniższym diagramie, wskaż kolor, który, obrazuje związek pomiędzy ceną samochodu Marysi, a jego wiekiem.  (Wskazówka: Odpowiedź zawiera nie tylko kolor diagramu, ale również odpowiednie równanie funkcji.)}
{zielony, \( p=12{,}0\cdot(0{,}88)^t \)}{żółty, \( p=14{,}5\cdot(1{,}15)^t \)}{pomarańczowy, \( p=12{,}0\cdot(1{,}12)^t \)}{niebieski, \( p=12{,}0\cdot(0{,}88)^t \)}{}{}}
\LANGcs{
\Question{
Rebeka si koupila nové auto za \( 12\,000 \) eur. Hodnota Rebečina auta se snižuje o  \( 12\  \% \) ročně. Zuzka si ve stejné době jako Rebeka koupila nové auto za \( 14\,500 \) eur. Hodnota Zuzčina auta se snižuje o \( 15\  \% \) ročně. Nechť \( p \) je hodnota auta v tisících eur a \( t \) je stáří auta v letech. Rozhodněte, jakou barvu na níže uvedeném obrázku má graf, který vyjadřuje vztah mezi hodnotou Rebečina auta a jeho stářím. (Poznámka: Odpovědi obsahují kromě barvy grafu i předpis příslušné funkce.)}
{zelená, \( p=12{,}0\cdot(0{,}88)^t \)}{žlutá, \( p=14{,}5\cdot(1{,}15)^t \)}{oranžová, \( p=12{,}0\cdot(1{,}12)^t \)}{modrá, \( p=12{,}0\cdot(0{,}88)^t \)}{}{}}
\LANGsk{
\Question{
Rebeka si kúpila nové auto za cenu \(12\,000\) eur. Depreciácia (pokles) hodnoty auta je \( 12\% \) ročne. Zuzka si tiež v tom istom čase ako Rebeka kúpila nové auto za cenu \( 14\,500 \) eur. Depreciácia (pokles) hodnoty Zuzkinho auta je \( 15\% \) ročne. Nech \( p \) je cena auta v tisícoch eur a \( t \) je vek auta v rokoch. Rozhodnite, akú farbu na nižšie uvedenom obrázku má graf, ktorý vyjadruje vzťah medzi cenou Rebekiného auta a jeho vekom. (Poznámka: Odpovede neobsahujú len farbu grafu, ale aj predpis funkcie.)}
{zelená, \( p=12{,}0\cdot(0{,}88)^t \)}{žltá, \( p=14{,}5\cdot(1{,}15)^t \)}{oranžová, \( p=12{,}0\cdot(1{,}12)^t \)}{modrá, \( p=12{,}0\cdot(0{,}88)^t \)}{}{}}
\LANGes{
\Question{
María compró un coche nuevo por el precio de \(12\,000\) euros y el valor del coche baja el \( 12\% \) al año. Susana compró un coche nuevo el mismo año que María y pagó \( 14\,500 \) euros. El valor del coche de Susana baja el \( 15\% \) al año. Sea \( p \) el precio de coche en miles de euros y \( t \) el tiempo del coche en años. En el gráfico de bajo, determina qué color de la gráfica representa la relación entre el precio del coche de María y su edad. (Ojo: Las respuestas incluyen no solo el color de la gráfica sino también la ecuación de la función adecuada.)}
{verde, \( p=12.0\cdot(0.88)^t \)}{amarillo, \( p=14.5\cdot(1.15)^t \)}{naranja, \( p=12.0\cdot(1.12)^t \)}{azul, \( p=12.0\cdot(0.88)^t \)}{}{}}
\End

\Data\ID{1103082503}
\Updated{2021-04-11 09:04:08}
\Level{C}
\SubArea{Exponential functions}
\IMG{1103082503.pdf}
\LANGen{
\Question{
At year \( 2000 \) (\(t=0\)) a small village has a population of \( 350 \) people. The graph in the picture below shows the population function for the next \( 50 \) years. One of the following statements is true. Which one is it?}
{The village population declines by \( 3.8\% \) per year.}{The village population grows by \( 3.8\% \) per year.}{The village population declines by \( 5 \) people per year.}{The village population grows by \( 5 \) people per year.}{}{}}
\LANGpl{
\Question{
W roku  \( 2000 \) (\(t=0\)) mała wioska liczy \( 350 \) osób. Wykres poniżej przedstawia funkcję dotyczącą populacji na kolejne \( 50 \) lat. Które z poniższych stwierdzeń jest zgodne z prawdą?}
{Populacja wioski zmniejsza się o  \( 3{,}8\% \) rocznie.}{Populacja wioski wzrasta o \( 3{,}8\% \) rocznie.}{Populacja wioski maleje o \( 5 \) osób rocznie.}{Populacja wzrasta o  \( 5 \) osób rocznie.}{}{}}
\LANGcs{
\Question{
V roce \( 2000 \) (\(t=0\)) byla populace malé vesnice \( 350 \) obyvatel. Graf na obrázku znázorňuje funkci populace v následujících \( 50 \) letech. Které z následujících tvrzení je pravdivé?}
{Pokles populace je \( 3{,}8\% \) ročně.}{Nárůst populace je \( 3{,}8\% \) ročně.}{Pokles populace je \( 5 \) obyvatel ročně.}{Nárůst populace je \( 5 \) obyvatel ročně.}{}{}}
\LANGsk{
\Question{
V roku \( 2000 \) (\(t=0\)) bola populácia malej dedinky \( 350 \) obyvateľov. Graf na obrázku znázorňuje funkciu populácie v nasledujúcich \( 50 \) rokoch. Ktoré z nasledujúcich tvrdení je pravdivé?}
{Pokles populácie je \( 3{,}8\% \) ročne.}{Nárast populácie je \( 3{,}8\% \) ročne.}{Pokles populácie je \( 5 \) obyvateľov ročne.}{Nárast populácie je \( 5 \) obyvateľov ročne.}{}{}}
\LANGes{
\Question{
En el año \( 2000 \) (\(t=0\)) un pueblo pequeño tenía \( 350 \) habitantes. La gráfica muestra la función de la población para los siguientes \( 50 \) años. ¿Cuál de los siguientes enunciados es verdadero?}
{La población disminuye un \( 3.8\% \) al año.}{La población crece un \( 3.8\% \) al año.}{La población disminuye en \( 5 \) personas cada año.}{La población crece en \( 5 \) personas cada año.}{}{}}
\End

\Data\ID{1103082704}
\Updated{2020-07-15 03:07:12}
\Level{C}
\SubArea{Exponential functions}
\IMG{1103082704.pdf}
\LANGen{
\Question{
Function \( f \) is given completely by the next graph. Identify which of the following statements is true.}
{\( f(x)=2^{|x|};\ x\in[-2;2] \)}{\( f(x)=\left|2^x\right|;\ x\in[-2;2] \)}{\( f(x)=\left|x^2+1\right|;\ x\in[-2;2] \)}{\( f(x)=\left|2^{-x}\right|;\ x\in[-2;2] \)}{}{}}
\LANGpl{
\Question{
Funkcja \( f \) przedstawiona została za pomocą poniższego wykresu. Które z poniższych stwierdzeń jest prawdziwe?}
{\( f(x)=2^{|x|};\ x\in\langle-2;2\rangle \)}{\( f(x)=\left|2^x\right|;\ x\in\langle-2;2\rangle \)}{\( f(x)=\left|x^2+1\right|;\ x\in\langle-2;2\rangle \)}{\( f(x)=\left|2^{-x}\right|;\ x\in\langle-2;2\rangle \)}{}{}}
\LANGcs{
\Question{
Funkce \( f \) je dána následujícím grafem. Určete, které z následujících tvrzení je pravdivé.}
{\( f(x)=2^{|x|};\ x\in\langle-2;2\rangle \)}{\( f(x)=\left|2^x\right|;\ x\in\langle-2;2\rangle \)}{\( f(x)=\left|x^2+1\right|;\ x\in\langle-2;2\rangle \)}{\( f(x)=\left|2^{-x}\right|;\ x\in\langle-2;2\rangle \)}{}{}}
\LANGsk{
\Question{
Funkcia \( f \) je daná nasledujúcim grafom. Vyberte správne tvrdenie o funkcii \( f \).}
{\( f(x)=2^{|x|};\ x\in\langle-2;2\rangle \)}{\( f(x)=\left|2^x\right|;\ x\in\langle-2;2\rangle \)}{\( f(x)=\left|x^2+1\right|;\ x\in\langle-2;2\rangle \)}{\( f(x)=\left|2^{-x}\right|;\ x\in\langle-2;2\rangle \)}{}{}}
\LANGes{
\Question{
La función \( f \) está definida por la siguiente gráfica. Identifica cuál de las siguientes afirmaciones es verdadera.}
{\( f(x)=2^{|x|};\ x\in[-2;2] \)}{\( f(x)=\left|2^x\right|;\ x\in[-2;2] \)}{\( f(x)=\left|x^2+1\right|;\ x\in[-2;2] \)}{\( f(x)=\left|2^{-x}\right|;\ x\in[-2;2] \)}{}{}}
\End

\Data\ID{9000003607}
\Updated{2020-07-15 03:07:34}
\Level{C}
\SubArea{Exponential functions}
\IMG{000036_07-figure0.pdf}
\LANGen{
\Question{
The function \(f(x) = \left (\frac{1}
{3}\right )^{x}\) 
is graphed in the picture. Identify a possible analytic expression for the function
\(g\).}
{\(y = 3^{|x|}- 1\)}{\(y = \left |\left (\frac{1}
{3}\right )^{x} - 1\right |\)}{\(y = \left (\frac{1}
{3}\right )^{|x|}- 1\)}{\(y = \left (\frac{1}
{3}\right )^{|x-1|}\)}{\(y = \left |3^{x} - 1\right |\)}{\(y = 3^{|x-1|}\)}}
\LANGpl{
\Question{
Funkcja  \(f(x) = \left (\frac{1}
{3}\right )^{x}\) 
jest przedstawiona za pomocą wykresu poniżej. Określ możliwy wzór funkcji
\(g\).}
{\(y = 3^{|x|}- 1\)}{\(y = \left |\left (\frac{1}
{3}\right )^{x} - 1\right |\)}{\(y = \left (\frac{1}
{3}\right )^{|x|}- 1\)}{\(y = \left (\frac{1}
{3}\right )^{|x-1|}\)}{\(y = \left |3^{x} - 1\right |\)}{\(y = 3^{|x-1|}\)}}
\LANGcs{
\Question{
Na obrázku jsou grafy funkcí \(f(x) = \left (\frac{1}
{3}\right )^{x}\) 
a \(g\). Jaký předpis
odpovídá funkci \(g\)?}
{\(y = 3^{|x|}- 1\)}{\(y = \left |\left (\frac{1}
{3}\right )^{x} - 1\right |\)}{\(y = \left (\frac{1}
{3}\right )^{|x|}- 1\)}{\(y = \left (\frac{1}
{3}\right )^{|x-1|}\)}{\(y = \left |3^{x} - 1\right |\)}{\(y = 3^{|x-1|}\)}}
\LANGsk{
\Question{
Na obrázku sú grafy funkcií \(f(x) = \left (\frac{1}
{3}\right )^{x}\) 
a \(g\). Aký predpis
zodpovedá funkcii \(g\)?}
{\(y = 3^{|x|}- 1\)}{\(y = \left |\left (\frac{1}
{3}\right )^{x} - 1\right |\)}{\(y = \left (\frac{1}
{3}\right )^{|x|}- 1\)}{\(y = \left (\frac{1}
{3}\right )^{|x-1|}\)}{\(y = \left |3^{x} - 1\right |\)}{\(y = 3^{|x-1|}\)}}
\LANGes{
\Question{
En la imagen están las funciones \(f(x) = \left (\frac{1}
{3}\right )^{x}\) y \(g\).
Identifica la expresión analítica para la función
\(g\).}
{\(y = 3^{|x|}- 1\)}{\(y = \left |\left (\frac{1}
{3}\right )^{x} - 1\right |\)}{\(y = \left (\frac{1}
{3}\right )^{|x|}- 1\)}{\(y = \left (\frac{1}
{3}\right )^{|x-1|}\)}{\(y = \left |3^{x} - 1\right |\)}{\(y = 3^{|x-1|}\)}}
\End
